% Options for packages loaded elsewhere
\PassOptionsToPackage{unicode}{hyperref}
\PassOptionsToPackage{hyphens}{url}
\documentclass[
  a4paper,
]{article}
\usepackage{xcolor}
\usepackage[margin=25mm]{geometry}
\usepackage{amsmath,amssymb}
\setcounter{secnumdepth}{-\maxdimen} % remove section numbering
\usepackage{iftex}
\ifPDFTeX
  \usepackage[T1]{fontenc}
  \usepackage[utf8]{inputenc}
  \usepackage{textcomp} % provide euro and other symbols
\else % if luatex or xetex
  \usepackage{unicode-math} % this also loads fontspec
  \defaultfontfeatures{Scale=MatchLowercase}
  \defaultfontfeatures[\rmfamily]{Ligatures=TeX,Scale=1}
\fi
\usepackage{lmodern}
\ifPDFTeX\else
  % xetex/luatex font selection
\fi
% Use upquote if available, for straight quotes in verbatim environments
\IfFileExists{upquote.sty}{\usepackage{upquote}}{}
\IfFileExists{microtype.sty}{% use microtype if available
  \usepackage[]{microtype}
  \UseMicrotypeSet[protrusion]{basicmath} % disable protrusion for tt fonts
}{}
\makeatletter
\@ifundefined{KOMAClassName}{% if non-KOMA class
  \IfFileExists{parskip.sty}{%
    \usepackage{parskip}
  }{% else
    \setlength{\parindent}{0pt}
    \setlength{\parskip}{6pt plus 2pt minus 1pt}}
}{% if KOMA class
  \KOMAoptions{parskip=half}}
\makeatother
\usepackage{color}
\usepackage{fancyvrb}
\newcommand{\VerbBar}{|}
\newcommand{\VERB}{\Verb[commandchars=\\\{\}]}
\DefineVerbatimEnvironment{Highlighting}{Verbatim}{commandchars=\\\{\}}
% Add ',fontsize=\small' for more characters per line
\newenvironment{Shaded}{}{}
\newcommand{\AlertTok}[1]{\textcolor[rgb]{1.00,0.00,0.00}{\textbf{#1}}}
\newcommand{\AnnotationTok}[1]{\textcolor[rgb]{0.38,0.63,0.69}{\textbf{\textit{#1}}}}
\newcommand{\AttributeTok}[1]{\textcolor[rgb]{0.49,0.56,0.16}{#1}}
\newcommand{\BaseNTok}[1]{\textcolor[rgb]{0.25,0.63,0.44}{#1}}
\newcommand{\BuiltInTok}[1]{\textcolor[rgb]{0.00,0.50,0.00}{#1}}
\newcommand{\CharTok}[1]{\textcolor[rgb]{0.25,0.44,0.63}{#1}}
\newcommand{\CommentTok}[1]{\textcolor[rgb]{0.38,0.63,0.69}{\textit{#1}}}
\newcommand{\CommentVarTok}[1]{\textcolor[rgb]{0.38,0.63,0.69}{\textbf{\textit{#1}}}}
\newcommand{\ConstantTok}[1]{\textcolor[rgb]{0.53,0.00,0.00}{#1}}
\newcommand{\ControlFlowTok}[1]{\textcolor[rgb]{0.00,0.44,0.13}{\textbf{#1}}}
\newcommand{\DataTypeTok}[1]{\textcolor[rgb]{0.56,0.13,0.00}{#1}}
\newcommand{\DecValTok}[1]{\textcolor[rgb]{0.25,0.63,0.44}{#1}}
\newcommand{\DocumentationTok}[1]{\textcolor[rgb]{0.73,0.13,0.13}{\textit{#1}}}
\newcommand{\ErrorTok}[1]{\textcolor[rgb]{1.00,0.00,0.00}{\textbf{#1}}}
\newcommand{\ExtensionTok}[1]{#1}
\newcommand{\FloatTok}[1]{\textcolor[rgb]{0.25,0.63,0.44}{#1}}
\newcommand{\FunctionTok}[1]{\textcolor[rgb]{0.02,0.16,0.49}{#1}}
\newcommand{\ImportTok}[1]{\textcolor[rgb]{0.00,0.50,0.00}{\textbf{#1}}}
\newcommand{\InformationTok}[1]{\textcolor[rgb]{0.38,0.63,0.69}{\textbf{\textit{#1}}}}
\newcommand{\KeywordTok}[1]{\textcolor[rgb]{0.00,0.44,0.13}{\textbf{#1}}}
\newcommand{\NormalTok}[1]{#1}
\newcommand{\OperatorTok}[1]{\textcolor[rgb]{0.40,0.40,0.40}{#1}}
\newcommand{\OtherTok}[1]{\textcolor[rgb]{0.00,0.44,0.13}{#1}}
\newcommand{\PreprocessorTok}[1]{\textcolor[rgb]{0.74,0.48,0.00}{#1}}
\newcommand{\RegionMarkerTok}[1]{#1}
\newcommand{\SpecialCharTok}[1]{\textcolor[rgb]{0.25,0.44,0.63}{#1}}
\newcommand{\SpecialStringTok}[1]{\textcolor[rgb]{0.73,0.40,0.53}{#1}}
\newcommand{\StringTok}[1]{\textcolor[rgb]{0.25,0.44,0.63}{#1}}
\newcommand{\VariableTok}[1]{\textcolor[rgb]{0.10,0.09,0.49}{#1}}
\newcommand{\VerbatimStringTok}[1]{\textcolor[rgb]{0.25,0.44,0.63}{#1}}
\newcommand{\WarningTok}[1]{\textcolor[rgb]{0.38,0.63,0.69}{\textbf{\textit{#1}}}}
\usepackage{longtable,booktabs,array}
\newcounter{none} % for unnumbered tables
\usepackage{calc} % for calculating minipage widths
% Correct order of tables after \paragraph or \subparagraph
\usepackage{etoolbox}
\makeatletter
\patchcmd\longtable{\par}{\if@noskipsec\mbox{}\fi\par}{}{}
\makeatother
% Allow footnotes in longtable head/foot
\IfFileExists{footnotehyper.sty}{\usepackage{footnotehyper}}{\usepackage{footnote}}
\makesavenoteenv{longtable}
\usepackage{graphicx}
\makeatletter
\newsavebox\pandoc@box
\newcommand*\pandocbounded[1]{% scales image to fit in text height/width
  \sbox\pandoc@box{#1}%
  \Gscale@div\@tempa{\textheight}{\dimexpr\ht\pandoc@box+\dp\pandoc@box\relax}%
  \Gscale@div\@tempb{\linewidth}{\wd\pandoc@box}%
  \ifdim\@tempb\p@<\@tempa\p@\let\@tempa\@tempb\fi% select the smaller of both
  \ifdim\@tempa\p@<\p@\scalebox{\@tempa}{\usebox\pandoc@box}%
  \else\usebox{\pandoc@box}%
  \fi%
}
% Set default figure placement to htbp
\def\fps@figure{htbp}
\makeatother
\setlength{\emergencystretch}{3em} % prevent overfull lines
\providecommand{\tightlist}{%
  \setlength{\itemsep}{0pt}\setlength{\parskip}{0pt}}
\usepackage{bookmark}
\IfFileExists{xurl.sty}{\usepackage{xurl}}{} % add URL line breaks if available
\urlstyle{same}
\hypersetup{
  pdftitle={The Fifth Thought Experiment: Was There No Hidden State? - A Self-Audit of the Two-State Orbit Generation of the Fourth Thought Experiment{,} Tracking the State Closure from Five States to the Sixth Persistent State \textbackslash mathcal N=\textbackslash nu -},
  pdfauthor={Noriaki Kihara (ORCID 0009-0004-6753-4020)},
  hidelinks,
  pdfcreator={LaTeX via pandoc}}

\title{The Fifth Thought Experiment: Was There No Hidden State? - A
Self-Audit of the Two-State Orbit Generation of the Fourth Thought
Experiment, Tracking the State Closure from Five States to the Sixth
Persistent State \(\mathcal N=\nu\) -}
\author{Noriaki Kihara (ORCID 0009-0004-6753-4020)}
\date{2026-09-26 / v2.1 / Version DOI 10.5281/zenodo.22973087}

\sloppy
\emergencystretch=3em
\begin{document}
\maketitle

\textbf{Author:} Noriaki Kihara\\
\textbf{ORCID:} 0009-0004-6753-4020\\
\textbf{Version DOI:} 10.5281/zenodo.22973087\\
\textbf{Concept DOI:} 10.5281/zenodo.22973086\\
\textbf{Version:} v2.1\\
\textbf{Date:} 2026-09-26

\begin{center}\rule{0.5\linewidth}{0.5pt}\end{center}

\subsection{Abstract}\label{abstract}

In the fourth thought experiment {[}4{]}, a state-dependent action
\(S_n\) was applied sequentially to the two-component internal state

\[
X_n=\begin{pmatrix}a_n\\b_n\end{pmatrix}
\]

and, from the quadratic readout

\[
\Phi(X_n)=
\begin{pmatrix}
a_n^2-b_n^2\\
2a_nb_n
\end{pmatrix}
\]

a family of gravitational orbits including the periastron advance and
the decay by the leading 2.5PN radiation reaction was reconstructed.
When that implementation is audited as a flow of information, however,
the processing phase \(\chi\), the semi-latus rectum \(p\), the
eccentricity \(e\), the azimuth \(\phi\) and the clock \(t\) are held
and updated outside \(X_n\), and \(X_n\) is a derived representation
built from them. The reading that ``the orbit was generated by two
states alone'' is therefore not exact.

In the first stage of the audit, the five persistent logical states that
survive across macro steps,

\[
(\chi,p,e,\phi,t)
\]

were made explicit and re-encoded as

\[
Z_5=(U,P,E,H,Q)
=(e^{i\chi},p,e,e^{i\phi},e^{t/\tau_0}).
\]

The RK4 sequential processing of Paper 4 could be rearranged into a
synchronous state machine in which the evaluation points,
\(k_1,\ldots,k_4\), the midpoints, the composite increments and an
11-phase one-hot phase are explicit machine/work states. In long
comparisons over 12 conditions, the difference from the reference update
was at most \(3.38\times10^{-16}\) in \(H\) and \(4.00\times10^{-15}\)
in \(Q\).

A further audit, however, found that the symmetric mass ratio \(\nu\) in
the radiation-reaction kernel

\[
C_{\rm RR}=\frac85\nu r^{-3}
\]

remained as an external argument outside the state. The system after the
first audit was therefore in fact

\[
Z_5'=F(Z_5;\nu)
\]

and was not fully state-closed. In this paper it is internalized as the
sixth persistent state,

\[
\boxed{
Z_6=(U,P,E,H,Q,\mathcal N),\qquad \mathcal N=\nu
}
\]

and held in all 11 phases as the identity row

\[
\mathcal N'=
\sum_{j=0}^{10}q_j(1\cdot\mathcal N)
=\mathcal N
\]

inside the same synchronous map. In the control experiments in which the
external \texttt{nu} argument was removed from the new transition
function, over 12 conditions \(\times\) 4000 macrosteps the first five
components of the old 5-state and the new 6-state systems, and the
stored values of \(r,x,y,t\), were bitwise identical, and the bitwise
mismatch of \(\mathcal N\) was also 0. Over 3 representative conditions
\(\times\) 12000 macrosteps the maximum difference of the compared
quantities was also 0, and there were 0 violations of the microstep
identity of \(\mathcal N\). The error profile against the stored Paper 4
C1 data also coincided exactly with the old system.

Furthermore, raising the angular series to \(N=12,16,20,24\) reduced the
azimuth error after three orbits for the most severe case
\(p=24.3,e=0.55\) to \(1.56\times10^{-7}\), \(9.03\times10^{-10}\),
\(5.54\times10^{-12}\) and \(3.55\times10^{-14}\). For the time
integration too, raising the order to 4, 6 and 8 was verified with the
same physical state kept. Therefore, at least in these experiments,
correcting the state closure and raising the numerical precision are
separate operations.

This paper does not claim that six states are mathematically minimal,
that the 11-phase machine is necessary in nature, that the whole
existing B scheme satisfies the strictest ``fully multiplicative
internal operation'' condition defined afterwards, or that general
relativity has been derived from this state system. What is shown is
that a self-audit of the fourth thought experiment corrects the state
closure from an apparent two states to five persistent states and
further to six persistent states including the external \(\nu\), and
that an orbit generation comparable with the known gravity model can be
reconstructed as explicit states, synchronous updates and a non-feedback
readout.

\textbf{Keywords:} hidden state, state closure, synchronous update,
multiplicative state, state-dependent action, symmetric mass ratio,
post-Newtonian approximation, discrete dynamics, numerical audit,
thought experiment

\begin{center}\rule{0.5\linewidth}{0.5pt}\end{center}

\section{0. Position of this paper}\label{position-of-this-paper}

This paper is the fifth thought experiment under the research project
design {[}0{]}.

The first thought experiment {[}1{]} examined, without first placing a
background spacetime or a metric, the minimal construction that
distinguishes inertial, uniformly accelerated and rotating motion, from
two degrees of freedom of undetermined meaning and the conservation of
an area readout.

The second thought experiment {[}2{]} applied a fixed action \(S\)
repeatedly to those two values and constructed inverse-square motion
from a product readout and an area clock.

The third thought experiment {[}3{]} reversed the direction and examined
whether the mass ratio, the relative distance and the relative time can
be read out from known Newtonian / PN two-body orbits.

The fourth thought experiment {[}4{]}, in order to bring the orbit
generation computed externally in the third thought experiment back to
an internal interaction, constructed the state-dependent action

\[
X_{n+1}=S_nX_n,
\qquad
S_n=\mathcal S(\mathcal R_n),
\qquad
\det S_n=1.
\]

The purpose of the fifth thought experiment is not to extend that result
further. The question is

\begin{center}
\fbox{\parbox{0.90\linewidth}{\centering Was the fourth thought experiment really closed by the states it made explicit alone?}}
\end{center}

In this paper the labels are used with the following meaning.

{\def\LTcaptype{none} % do not increment counter
\begin{longtable}[]{@{}
  >{\raggedright\arraybackslash}p{(\linewidth - 2\tabcolsep) * \real{0.5000}}
  >{\raggedright\arraybackslash}p{(\linewidth - 2\tabcolsep) * \real{0.5000}}@{}}
\toprule\noalign{}
\begin{minipage}[b]{\linewidth}\raggedright
Label
\end{minipage} & \begin{minipage}[b]{\linewidth}\raggedright
Meaning
\end{minipage} \\
\midrule\noalign{}
\endhead
\bottomrule\noalign{}
\endlastfoot
{[}Principle{]} & Computational principle maintained in this paper for
the self-audit \\
{[}Design assumption{]} & Assumption made for the present representation
and implementation; may be removed later \\
{[}Construction{]} & Procedure that constructs an action from known
formulas or explicit states \\
{[}Numerical{]} & Result confirmed by numerical comparison with
programs \\
{[}Cross-check{]} & Comparison with known physics or existing
mathematics; not a derivation from first principles \\
{[}Unproven{]} & Matter that cannot be decided by the present
experiments \\
\end{longtable}
}

The most important policy of this paper is \textbf{not to presuppose the
success of the fourth thought experiment}. If a hidden state is found,
it is not conveniently ignored but moved into the state. Indeed, in this
paper, after making the five persistent states explicit in the first
audit, the external \(\nu\) was found in addition and internalized as
the sixth persistent state. This process of correction itself is
recorded as a result.

\section{1. Motivation: was the orbit really generated by two states
alone?}\label{motivation-was-the-orbit-really-generated-by-two-states-alone}

The surface state of the fourth thought experiment is

\[
X_n=(a_n,b_n)^{\mathsf T}
\]

and the orbit was read out from

\[
Y_n=\Phi(X_n).
\]

In the actual program, however, at each step

\[
\chi_n,
\quad p_n,
\quad e_n,
\quad \phi_n,
\quad t_n
\]

were held separately, from which

\[
r_n=\frac{p_n}{1+e_n\cos\chi_n},
\qquad
\theta_n=\frac{\phi_n}{2}
\]

were obtained, and then \(S_n\) and \(X_n\) were constructed {[}4{]}.
The actual flow of information was therefore

\[
(\chi,p,e,\phi,t)
\longrightarrow
(r,\theta)
\longrightarrow
S_n
\longrightarrow
X
\longrightarrow
Y
\]

and \(X\) was a derived representation obtained from the upstream
states.

The first audit made this upstream information explicit as five
persistent states. But when the state dependence was tracked further
into the radiation-reaction kernel, the symmetric mass ratio \(\nu\) was
being passed from outside the state. The path of the state audit is
therefore

\[
\boxed{
X=(a,b)
\;\longrightarrow\;
Z_5=(U,P,E,H,Q)
\;\longrightarrow\;
Z_6=(U,P,E,H,Q,\mathcal N)
}
\]

\begin{figure}
\centering
\pandocbounded{\includegraphics[keepaspectratio,alt={Figure 1. The path of the state audit: 2 to 5 to 6}]{figs/figure01.png}}
\caption{Figure 1. The path of the state audit: 2 to 5 to 6}
\end{figure}

\textbf{Figure 1.} Tracing the apparent two states of Paper 4 upstream
reveals five persistent states, and the additional audit found the
external \(\nu\). After the revision, \(\mathcal N=\nu\) is included as
the sixth persistent state in the same state map. The RK4 work states
and the 11-phase one-hot phase \(q\) are separate machine/work states;
this is not a claim that ``the total number of real scalar degrees of
freedom is 6''.

So as not to lose sight of the object of the audit itself, Figure 2
shows the orbit of the representative condition C1 that Paper 4
regenerated from the sequential interaction. This paper is not a paper
deriving this orbit from a new law of gravity; it is a paper that
self-audits which persistent information this existing orbit generation
depended on. In Figure 2, the stored Paper 4 C1 orbit of 12001 points is
overlaid with the orbit of 12000 macrosteps (3 radial periods)
regenerated independently by the generator in which the external \(\nu\)
is internalized as the sixth state \(\mathcal N\). During the
generation, the stored C1 data are not fed back into the state update.

\begin{figure}
\centering
\pandocbounded{\includegraphics[keepaspectratio,alt={Figure 2. The Paper 4 C1 orbit and the orbit regenerated after internalizing the sixth state}]{figs/figure_orbit_c1_paper4_vs_6state.png}}
\caption{Figure 2. The Paper 4 C1 orbit and the orbit regenerated after
internalizing the sixth state}
\end{figure}

\textbf{Figure 2.} The stored C1 orbit of Paper 4 and the orbit
regenerated by \(Z_6=(U,P,E,H,Q,\mathcal N)\). The 12000 macrosteps
correspond to 3 radial periods and include the periastron advance and
the orbital decay by the radiation reaction. The two curves overlap at
the drawing resolution. The maximum absolute differences against the
stored C1 are \(|\Delta x|=3.85\times10^{-12}\),
\(|\Delta y|=5.11\times10^{-12}\), \(|\Delta r|=6.89\times10^{-13}\) and
\(|\Delta t|=1.94\times10^{-10}\). This is not an independent
verification against a new physical model; it visualizes that the orbit
generation of Paper 4 can be kept as the same output after internalizing
the sixth state.

The question of this paper is divided into four.

\begin{center}
\fbox{\parbox{0.90\linewidth}{\centering (1) When all the information that determines the next state is counted as state, what remains?}}
\end{center}

\begin{center}
\fbox{\parbox{0.90\linewidth}{\centering (2) Does any case-dependent physical parameter flow in from outside the state?}}
\end{center}

\begin{center}
\fbox{\parbox{0.90\linewidth}{\centering (3) Can the intermediate update values and the processing stages be written as a synchronous map without hiding them?}}
\end{center}

\begin{center}
\fbox{\parbox{0.90\linewidth}{\centering (4) Can raising the approximation order be separated from adding a physical state?}}
\end{center}

\section{2. Purpose and scope of the
claims}\label{purpose-and-scope-of-the-claims}

The purpose of this paper is not to derive a new exact solution of
general relativity. Nor does it claim that the state transition rule
constructed here is a fundamental law of nature.

The purpose is to audit the external clock, the past states, the
processing step number, the intermediate update values, the
case-dependent parameters and the implicit registers that ordinary
numerical computation tends to use implicitly, and to examine how far
the computation can be closed into the form

\begin{center}
\fbox{\parbox{0.90\linewidth}{\centering present explicit state$\rightarrow$interaction$\rightarrow$next explicit state$\rightarrow$non-feedback readout}}
\end{center}

\subsubsection{What was confirmed in this
paper}\label{what-was-confirmed-in-this-paper}

\begin{enumerate}
\def\labelenumi{\arabic{enumi}.}
\tightlist
\item
  The implementation was not closed by the two states \(X=(a,b)\) of
  Paper 4 alone.
\item
  In the first audit, making the five quantities \(\chi,p,e,\phi,t\)
  explicit as persistent states allowed the orbit computation of Paper 4
  to be reconstructed.
\item
  However, \(\nu\) remained as an external argument in the
  radiation-reaction kernel, and the five-state system was not
  completely closed.
\item
  Internalizing \(\mathcal N=\nu\) as the sixth persistent state and
  implementing \(\mathcal N'=\mathcal N\) as an identity row of the same
  one-hot sum-of-products map removed the external \texttt{nu} argument
  without changing the old five states, the orbit or the clock bitwise.
\item
  The RK4 sequential computation of Paper 4 could be reconstructed as an
  11-phase synchronous state machine with explicit processing stages and
  work registers.
\item
  Raising the order of the angular series and of the numerical
  integration is independent of adding the sixth state, and no new
  persistent physical state had to be introduced to improve the
  precision.
\end{enumerate}

\subsubsection{What is not proved in this
paper}\label{what-is-not-proved-in-this-paper}

\begin{enumerate}
\def\labelenumi{\arabic{enumi}.}
\tightlist
\item
  That six persistent logical states are mathematically minimal.
\item
  That the 11-phase machine state is necessary in nature.
\item
  That a single fixed action \(S\) is impossible in principle and a
  sequence of actions is always necessary.
\item
  That the whole existing B scheme satisfies the strictest ``fully
  multiplicative internal operation'' condition defined afterwards.
\item
  That the limit of raising the approximation order to infinity
  coincides with the exact solution of general relativity.
\item
  That the interaction rule of this paper is a fundamental law of
  nature.
\end{enumerate}

\section{3. Rules of the self-audit}\label{rules-of-the-self-audit}

\subsection{3.1 Make the next state from the present state
alone}\label{make-the-next-state-from-the-present-state-alone}

Let the set of states be

\[
\Psi_n=(\psi_{1,n},\ldots,\psi_{N,n}).
\]

In one interaction, the present state is read once and the next state is
written once.

\[
\boxed{
\psi_{i,n+1}=F_i(\Psi_n)
}
\]

must hold.

Using a \(\psi_{j,n+1}\) computed earlier within the same single
interaction for the computation of another component is forbidden.

\begin{center}\rule{0.5\linewidth}{0.5pt}\end{center}

\subsection{3.2 Distinguish incremental assignment from sum of
products}\label{distinguish-incremental-assignment-from-sum-of-products}

In a program,

\begin{Shaded}
\begin{Highlighting}[]
\NormalTok{psi += delta}
\end{Highlighting}
\end{Shaded}

can be read as a sequential update that implicitly holds the old value
and the new value, so it is not adopted as a fundamental interaction.

On the other hand, a \textbf{sum of products} computed at once from the
present state alone,

\[
\psi_i'
=
\sum_{\alpha}
c_{i\alpha}
\prod_j \psi_j^{m_{i\alpha j}},
\]

and written at once, is permitted.

For example,

\[
P'=P+D_p
\]

does not violate the rule of this paper as long as it is computed at
once, not as a sequential assignment but as the sum of products of two
present states,

\[
P'=1\cdot P+1\cdot D_p.
\]

\begin{center}\rule{0.5\linewidth}{0.5pt}\end{center}

\subsection{3.3 Do not refer to the number of operations from
outside}\label{do-not-refer-to-the-number-of-operations-from-outside}

The action must not be selected by an external loop counter.

If the processing stages have to be distinguished, a one-hot phase state

\[
q_n=(q_{0,n},\ldots,q_{m-1,n})
\]

is placed inside the system and updated by the cyclic permutation

\[
\boxed{
q_{n+1}=C_m q_n,
\qquad C_m^m=I.
}
\]

The selection of the action is also written, in the formulas, as the sum
of products

\[
\boxed{
\Psi_{n+1}
=
\sum_{j=0}^{m-1} q_{j,n}F_j(\Psi_n).
}
\]

In the actual audit code, since \(q\) is exactly one-hot,
\texttt{argmax(q)} and a branch are used as an acceleration, but this is
treated not as physical external information but as an implementation of
the one-hot selection above. Therefore, \textbf{the fundamental
definition in the formulas is the sum-of-products selection by \(q\)},
and \texttt{argmax} itself is not regarded as a fundamental interaction.

\begin{center}\rule{0.5\linewidth}{0.5pt}\end{center}

\subsection{3.4 The readout is not written back into the
state}\label{the-readout-is-not-written-back-into-the-state}

A map for observation,

\[
O=\mathcal O(\Psi),
\]

is permitted.

However, the procedure of returning to additive coordinates, in
particular using

\[
\log,
\]

performing the conventional additive computation and exponentiating
again,

\[
\Psi
\rightarrow
\log\Psi
\rightarrow
\text{additive update}
\rightarrow
\exp
\rightarrow
\Psi',
\]

is not accepted as a fundamental model.

The readout is one-directional:

\[
\boxed{
\text{state}\rightarrow\text{observation}.
}
\]

\begin{center}\rule{0.5\linewidth}{0.5pt}\end{center}

\subsection{3.5 Constants}\label{constants}

A sum of products with fixed constants is permitted if their origin is
clear.

For example, the RK4 coefficients

\[
\frac16(1,2,2,1)
\]

and the coefficients of the series

\[
c_{m+1}=\frac{2m+1}{m+1}c_m
\]

have a clear derivation.

On the other hand, free fitting coefficients used only to approach a
known orbit are not used.

\begin{figure}
\centering
\pandocbounded{\includegraphics[keepaspectratio,alt={Figure 3. State closure and synchronous update in the revised version}]{figs/figure02.png}}
\caption{Figure 3. State closure and synchronous update in the revised
version}
\end{figure}

\textbf{Figure 3.} The scope of the audit after the revision. The old
full state is read at once, and the next full state is generated by the
same microstep map. The readouts \(x,y,r,t\) are not returned to the
state update. The identity holding of \(\mathcal N\) is also included in
the same map. On the other hand, the minimality of the six states, the
necessity of the 11-phase machine in nature, and the full
multiplicativity of the whole existing B scheme are not claimed.

\begin{center}\rule{0.5\linewidth}{0.5pt}\end{center}

\section{4. Position among existing
theories}\label{position-among-existing-theories}

This paper does not claim the individual computational concepts
themselves as new mathematics.

That in an effective description projected onto a small number of
variables the influence of the eliminated degrees of freedom can appear
as a memory term is represented by the projection formalism of Zwanzig
{[}5{]}. This paper does not use that result to derive the six states;
it refers to it only as the background for \textbf{auditing where the
missing information resides when the system does not close with the
present state alone}.

For representing the evolution as a product of several actions, the
product formula of Trotter {[}6{]} gives a classical precedent. The
11-phase state machine of this paper is not the Trotter decomposition
itself, but treating several actions as an explicit sequence is itself
within existing mathematics.

For numerical integration, the systematic construction of high-order
Runge--Kutta processes has long been studied by Butcher {[}7{]}. The
purpose of examining the raising of the order to 4, 6 and 8 in this
paper is not to propose a new integration method but \textbf{to
distinguish the depth of the numerical approximation from the number of
persistent physical states}.

On the gravity side, the Darwin parameterization {[}8{]} is used for the
angular progression of bound Schwarzschild orbits, and the
osculating-orbit formalism {[}9{]} is referred to as the closely related
standard background for the construction of orbital elements that evolve
in time. The general background of the post-Newtonian approximation and
the radiation reaction is placed in the review by Blanchet {[}10{]}, and
the lineage of the radiation-reaction expressions including the leading
2.5PN radiation-reaction term in Blanchet, Faye and Trestini {[}11{]}.

What is specific to this paper is therefore not a new invention of these
existing elements. It is to have audited, for the implementation of the
fourth thought experiment, the state leakage, the intermediate update
values, the processing phases and the feedback of the readout at the
same time, and to have corrected the state closure including the
external \(\nu\) found on the way.

\section{5. State audit of Paper 4}\label{state-audit-of-paper-4}

\subsection{\texorpdfstring{5.1 The surface state \(X=(a,b)\) was not
the fundamental
state}{5.1 The surface state X=(a,b) was not the fundamental state}}\label{the-surface-state-xab-was-not-the-fundamental-state}

Paper 4 displayed

\[
X_{n+1}=S_nX_n
\]

as the main state transition {[}4{]}.

In the implementation, however,

\[
(\chi_n,p_n,e_n,\phi_n,t_n)
\]

were updated first, and from

\[
r_n=\frac{p_n}{1+e_n\cos\chi_n},
\qquad
\theta_n=\frac{\phi_n}{2}
\]

the state

\[
X_n
=
\sqrt{r_n}
\begin{pmatrix}
\cos\theta_n\\
\sin\theta_n
\end{pmatrix}
\]

was constructed. \(X_n\) was therefore a compressed derived
representation of the upstream states, and by itself cannot determine
the next state.

\begin{center}\rule{0.5\linewidth}{0.5pt}\end{center}

\subsection{5.2 First audit: five persistent
states}\label{first-audit-five-persistent-states}

In the first audit, the states needed across macro steps in the present
representation were made explicit as

\[
(\chi,p,e,\phi,t):
\]

\begin{itemize}
\tightlist
\item
  \(\chi\): the processing phase,
\item
  \(p\): the semi-latus rectum,
\item
  \(e\): the eccentricity,
\item
  \(\phi\): the orbital azimuth,
\item
  \(t\): the accumulated coordinate of the clock readout.
\end{itemize}

These were re-encoded as

\[
Z_5=(U,P,E,H,Q)
=(e^{i\chi},p,e,e^{i\phi},e^{t/\tau_0}).
\]

Comparing the five-state synchronous map with the original computation
of Paper 4 over 12 conditions, the existing computation could be
reproduced including the orbit readout. At this stage it was once judged
that ``the system closed with five states''.

This conclusion, however, is corrected by the additional audit.

\subsubsection{Table 1. Final summary of the state
audit}\label{table-1.-final-summary-of-the-state-audit}

{\def\LTcaptype{none} % do not increment counter
\begin{longtable}[]{@{}
  >{\raggedright\arraybackslash}p{(\linewidth - 8\tabcolsep) * \real{0.2000}}
  >{\raggedright\arraybackslash}p{(\linewidth - 8\tabcolsep) * \real{0.2000}}
  >{\raggedright\arraybackslash}p{(\linewidth - 8\tabcolsep) * \real{0.2000}}
  >{\raggedright\arraybackslash}p{(\linewidth - 8\tabcolsep) * \real{0.2000}}
  >{\raggedright\arraybackslash}p{(\linewidth - 8\tabcolsep) * \real{0.2000}}@{}}
\toprule\noalign{}
\begin{minipage}[b]{\linewidth}\raggedright
Quantity
\end{minipage} & \begin{minipage}[b]{\linewidth}\raggedright
Role
\end{minipage} & \begin{minipage}[b]{\linewidth}\raggedright
Classification
\end{minipage} & \begin{minipage}[b]{\linewidth}\raggedright
Reason of the audit
\end{minipage} & \begin{minipage}[b]{\linewidth}\raggedright
Audit stage
\end{minipage} \\
\midrule\noalign{}
\endhead
\bottomrule\noalign{}
\endlastfoot
\(X=(a,b)\) & The two components displayed in Paper 4 & Derived
representation & Constructed from \(r,\phi\); does not close the next
state by itself & Paper 4 to first audit \\
\(U=e^{i\chi}\) & Processing phase & Persistent logical state & Holds
the periodic information of \(\chi\) & 5-state audit \\
\(P=p\) & Orbital element & Persistent logical state & Updated by the
radiation reaction & 5-state audit \\
\(E=e\) & Orbital element & Persistent logical state & Updated by the
radiation reaction & 5-state audit \\
\(H=e^{i\phi}\) & Orbital azimuth phase & Persistent logical state &
Holds the angular progression multiplicatively & 5-state audit \\
\(Q=e^{t/\tau_0}\) & Clock state & Persistent logical state & \(t\) is
obtained by \(\log\) only at readout & 5-state audit \\
\(\mathcal N=\nu\) & Symmetric mass ratio & Persistent logical state &
External argument in the old implementation; judged a state leakage in
the additional audit & Additional \(\nu\) audit \\
\(q_0,\ldots,q_{10}\) & 11-phase one-hot processing phase & Machine
state & Substitute for an external loop counter & 11-phase synchronous
audit \\
RK work registers & \(k_1,\ldots,k_4\) etc. & Work/machine state & Made
explicit so that intermediate update values are not held implicitly &
11-phase synchronous audit \\
\end{longtable}
}

\begin{center}\rule{0.5\linewidth}{0.5pt}\end{center}

\subsection{\texorpdfstring{5.3 Additional audit: the external \(\nu\)
was the sixth
state}{5.3 Additional audit: the external \textbackslash nu was the sixth state}}\label{additional-audit-the-external-nu-was-the-sixth-state}

Since the radiation-reaction kernel uses

\[
C_{\rm RR}=\frac85\nu r^{-3},
\]

the implementation after the first audit was in fact

\[
\boxed{Z_5'=F(Z_5;\nu)}.
\]

\(\nu\) is physical information given per case of fixed mass ratio, and
was supplied from outside the state map. Read strictly as a state
closure, therefore, the five-state system is not closed.

Accordingly,

\[
\boxed{
Z_6=(U,P,E,H,Q,\mathcal N),
\qquad
\mathcal N=\nu
}
\]

was taken, and the radiation-reaction kernel was replaced by

\[
C_{\rm RR}=\frac85\mathcal N r^{-3}.
\]

\(\mathcal N\) is held as an identity row with coefficient 1 in all 11
phases,

\[
\boxed{
\mathcal N_{m+1}
=
\sum_{j=0}^{10}q_{j,m}(1\cdot\mathcal N_m)
=
\mathcal N_m.
}
\]

This was implemented not as a copy \texttt{N\_next=N\_old} outside the
state machine but as the sixth row of the one-hot sum-of-products map.

The old transitions

\begin{Shaded}
\begin{Highlighting}[]
\NormalTok{micro\_fast(s, nu)}
\NormalTok{macro\_fast(s, nu)}
\end{Highlighting}
\end{Shaded}

become

\begin{Shaded}
\begin{Highlighting}[]
\NormalTok{micro\_fast6(s)}
\NormalTok{macro\_fast6(s)}
\end{Highlighting}
\end{Shaded}

and the new transition functions have no external \texttt{nu} argument.

\subsubsection{\texorpdfstring{Control results of the internalization of
\(\nu\)}{Control results of the internalization of \textbackslash nu}}\label{control-results-of-the-internalization-of-nu}

Over the 48,012 stored rows of 12 conditions \(\times\) 4000 macrosteps,
for the first five components of the old 5-state and the new 6-state
systems,

\[
\boxed{\Delta U=\Delta P=\Delta E=\Delta H=\Delta Q=0},
\]

and for the readouts too,

\[
\boxed{\Delta r=\Delta x=\Delta y=\Delta t=0}.
\]

In the comparison of the stored bit patterns in complex128 / float64 as
well, the bitwise mismatch between the old 5 states and the first 5
states of the new system was 0, and the bitwise mismatch between
\(\mathcal N\) and the initial \(\nu_0\) was also 0.

Furthermore, over 3 representative conditions \(\times\) 12000
macrosteps,

\[
\boxed{
\max|\Delta(U,P,E,H,Q,r,x,y,t)|=0
}
\]

and there were 0 violations of the microstep identity holding of
\(\mathcal N\). The error profiles of \(p,e,t,x,y\) against the stored
Paper 4 C1 CSV were also completely identical between the old 5-state
and the new 6-state systems. {[}Numerical{]}

What this control experiment shows is that \textbf{the physical formulas
were not changed and only the place where the information is kept was
moved from an external argument to the sixth state}. It is not the
result of adding a charged interaction or another physical degree of
freedom.

\begin{center}\rule{0.5\linewidth}{0.5pt}\end{center}

\subsection{5.4 The meaning of ``six
states''}\label{the-meaning-of-six-states}

The exact claim after the revision is

\begin{center}
\fbox{\parbox{0.90\linewidth}{\centering In the present representation, there are six persistent logical states held across macro steps.}}
\end{center}

This does not mean that the real scalar components of the whole state
machine number six. The \(k_1,\ldots,k_4\), the evaluation points, the
midpoints, the composite increments and the 11-phase one-hot phase
needed to implement RK4 without hidden sequential registers exist as
separate explicit machine/work states.

Also, that six is the mathematical minimum is unproven. Therefore

\begin{center}
\fbox{\parbox{0.90\linewidth}{\centering a self-audit of the present Paper 4 representation required at least six persistent logical states to be made explicit}}
\end{center}

is the conclusion of this paper.

\section{6. Re-encoding into multiplicative
states}\label{re-encoding-into-multiplicative-states}

To the additive coordinates obtained in the first audit,

\[
(\chi,p,e,\phi,t),
\]

the symmetric mass ratio \(\nu\) found in the additional audit is added,
and the persistent logical states after the revision are taken to be

\[
\boxed{
Z_6=(U,P,E,H,Q,\mathcal N).
}
\]

The definitions are

\[
U=e^{i\chi},
\qquad P=p,
\qquad E=e,
\qquad H=e^{i\phi},
\qquad Q=e^{t/\tau_0},
\qquad \mathcal N=\nu.
\]

The phases and the clock become the multiplicative updates

\[
U'=Ue^{i\Delta\chi},
\]

\[
H'=He^{i\Delta\phi},
\]

\[
Q'=Qe^{\Delta t/\tau_0}.
\]

In the fixed-mass-ratio model \(\mathcal N\) satisfies

\[
\mathcal N'=\mathcal N
\]

and is held as an identity row of the same state map.

On the other hand, since \(P,E\) appear directly as positive orbital
elements in the present physical formulas, this paper does not force
\(\log P,\log E\) to be persistent states.

{[}Design assumption{]} The present verification programs implement
\(U,H\) concisely as complex numbers. This is not, however, a conclusion
that the fundamental states must essentially be complex numbers. For
example, \(U=(\cos\chi,\sin\chi)\) can be taken as two real components
and updated by a rotation matrix. Therefore ``six states'' is a
statement about \textbf{the number of logical state quantities}, and
does not fix the minimal number of real scalar components.

\section{7. Scheme B: the 11-phase synchronous internal state
machine}\label{scheme-b-the-11-phase-synchronous-internal-state-machine}

\subsection{7.1 Why internal phases are
needed}\label{why-internal-phases-are-needed}

The radiation-reaction update of Paper 4 is RK4, and ordinary code
processes it sequentially as

\[
k_1\rightarrow y_2\rightarrow k_2\rightarrow y_3\rightarrow k_3\rightarrow y_4\rightarrow k_4.
\]

If this is regarded as one physical interaction, the values in the
middle of the processing become hidden registers.

Each intermediate value was therefore made an explicit state, and the
processing was decomposed into 11 microsteps.

The internal state is written schematically as

\[
\mathcal Z
=(Z,W,q),
\]

where \(Z_6=(U,P,E,H,Q,\mathcal N)\) is the persistent state, \(W\) is
the work state consisting of \(k_1,\ldots,k_4\), the evaluation points,
the composite increments, the midpoints and so on, and \(q\) is the
11-phase one-hot state.

One microstep is

\[
\boxed{
\mathcal Z_{m+1}
=
\sum_{j=0}^{10}q_{j,m}F_j(\mathcal Z_m)
}
\]

and at the same time

\[
\boxed{
q_{m+1}=C_{11}q_m.
}
\]

After 11 microsteps,

\[
q_{m+11}=q_m,
\]

and one macro state update is completed.

\begin{center}\rule{0.5\linewidth}{0.5pt}\end{center}

\subsection{7.2 The actual audit code}\label{the-actual-audit-code}

The following is the core of the control code after the internalization
of \(\nu\). \texttt{o} is the full state before the update and
\texttt{n} is the write destination of the next state; every branch
reads only the OLD state. The radiation reaction reads
\texttt{o{[}\textquotesingle{}N\textquotesingle{}{]}}, not an external
\texttt{nu}.

\begin{Shaded}
\begin{Highlighting}[]
\KeywordTok{def}\NormalTok{ micro\_fast6(s):}
\NormalTok{    j }\OperatorTok{=} \BuiltInTok{int}\NormalTok{(np.argmax(s[}\StringTok{\textquotesingle{}q\textquotesingle{}}\NormalTok{]))}
\NormalTok{    o }\OperatorTok{=}\NormalTok{ clone\_state(s)          }\CommentTok{\# OLD full state}
\NormalTok{    n }\OperatorTok{=}\NormalTok{ clone\_state(s)          }\CommentTok{\# next{-}state buffer}

\NormalTok{    U, P, E, N }\OperatorTok{=}\NormalTok{ o[}\StringTok{\textquotesingle{}U\textquotesingle{}}\NormalTok{], o[}\StringTok{\textquotesingle{}P\textquotesingle{}}\NormalTok{], o[}\StringTok{\textquotesingle{}E\textquotesingle{}}\NormalTok{], o[}\StringTok{\textquotesingle{}N\textquotesingle{}}\NormalTok{]}

    \ControlFlowTok{if}\NormalTok{ j }\OperatorTok{==} \DecValTok{0}\NormalTok{:}
\NormalTok{        n[}\StringTok{\textquotesingle{}k1\textquotesingle{}}\NormalTok{] }\OperatorTok{=}\NormalTok{ rr(U, P, E, N)}
    \ControlFlowTok{elif}\NormalTok{ j }\OperatorTok{==} \DecValTok{1}\NormalTok{:}
\NormalTok{        n[}\StringTok{\textquotesingle{}Us\textquotesingle{}}\NormalTok{] }\OperatorTok{=}\NormalTok{ U }\OperatorTok{*}\NormalTok{ PH2}
\NormalTok{        n[}\StringTok{\textquotesingle{}Ps\textquotesingle{}}\NormalTok{] }\OperatorTok{=}\NormalTok{ P }\OperatorTok{+} \FloatTok{0.5} \OperatorTok{*}\NormalTok{ H }\OperatorTok{*}\NormalTok{ o[}\StringTok{\textquotesingle{}k1\textquotesingle{}}\NormalTok{][}\DecValTok{0}\NormalTok{]}
\NormalTok{        n[}\StringTok{\textquotesingle{}Es\textquotesingle{}}\NormalTok{] }\OperatorTok{=}\NormalTok{ E }\OperatorTok{+} \FloatTok{0.5} \OperatorTok{*}\NormalTok{ H }\OperatorTok{*}\NormalTok{ o[}\StringTok{\textquotesingle{}k1\textquotesingle{}}\NormalTok{][}\DecValTok{1}\NormalTok{]}
    \ControlFlowTok{elif}\NormalTok{ j }\OperatorTok{==} \DecValTok{2}\NormalTok{:}
\NormalTok{        n[}\StringTok{\textquotesingle{}k2\textquotesingle{}}\NormalTok{] }\OperatorTok{=}\NormalTok{ rr(o[}\StringTok{\textquotesingle{}Us\textquotesingle{}}\NormalTok{], o[}\StringTok{\textquotesingle{}Ps\textquotesingle{}}\NormalTok{], o[}\StringTok{\textquotesingle{}Es\textquotesingle{}}\NormalTok{], N)}
    \ControlFlowTok{elif}\NormalTok{ j }\OperatorTok{==} \DecValTok{3}\NormalTok{:}
\NormalTok{        n[}\StringTok{\textquotesingle{}Us\textquotesingle{}}\NormalTok{] }\OperatorTok{=}\NormalTok{ U }\OperatorTok{*}\NormalTok{ PH2}
\NormalTok{        n[}\StringTok{\textquotesingle{}Ps\textquotesingle{}}\NormalTok{] }\OperatorTok{=}\NormalTok{ P }\OperatorTok{+} \FloatTok{0.5} \OperatorTok{*}\NormalTok{ H }\OperatorTok{*}\NormalTok{ o[}\StringTok{\textquotesingle{}k2\textquotesingle{}}\NormalTok{][}\DecValTok{0}\NormalTok{]}
\NormalTok{        n[}\StringTok{\textquotesingle{}Es\textquotesingle{}}\NormalTok{] }\OperatorTok{=}\NormalTok{ E }\OperatorTok{+} \FloatTok{0.5} \OperatorTok{*}\NormalTok{ H }\OperatorTok{*}\NormalTok{ o[}\StringTok{\textquotesingle{}k2\textquotesingle{}}\NormalTok{][}\DecValTok{1}\NormalTok{]}
    \ControlFlowTok{elif}\NormalTok{ j }\OperatorTok{==} \DecValTok{4}\NormalTok{:}
\NormalTok{        n[}\StringTok{\textquotesingle{}k3\textquotesingle{}}\NormalTok{] }\OperatorTok{=}\NormalTok{ rr(o[}\StringTok{\textquotesingle{}Us\textquotesingle{}}\NormalTok{], o[}\StringTok{\textquotesingle{}Ps\textquotesingle{}}\NormalTok{], o[}\StringTok{\textquotesingle{}Es\textquotesingle{}}\NormalTok{], N)}
    \ControlFlowTok{elif}\NormalTok{ j }\OperatorTok{==} \DecValTok{5}\NormalTok{:}
\NormalTok{        n[}\StringTok{\textquotesingle{}Us\textquotesingle{}}\NormalTok{] }\OperatorTok{=}\NormalTok{ U }\OperatorTok{*}\NormalTok{ PH}
\NormalTok{        n[}\StringTok{\textquotesingle{}Ps\textquotesingle{}}\NormalTok{] }\OperatorTok{=}\NormalTok{ P }\OperatorTok{+}\NormalTok{ H }\OperatorTok{*}\NormalTok{ o[}\StringTok{\textquotesingle{}k3\textquotesingle{}}\NormalTok{][}\DecValTok{0}\NormalTok{]}
\NormalTok{        n[}\StringTok{\textquotesingle{}Es\textquotesingle{}}\NormalTok{] }\OperatorTok{=}\NormalTok{ E }\OperatorTok{+}\NormalTok{ H }\OperatorTok{*}\NormalTok{ o[}\StringTok{\textquotesingle{}k3\textquotesingle{}}\NormalTok{][}\DecValTok{1}\NormalTok{]}
    \ControlFlowTok{elif}\NormalTok{ j }\OperatorTok{==} \DecValTok{6}\NormalTok{:}
\NormalTok{        n[}\StringTok{\textquotesingle{}k4\textquotesingle{}}\NormalTok{] }\OperatorTok{=}\NormalTok{ rr(o[}\StringTok{\textquotesingle{}Us\textquotesingle{}}\NormalTok{], o[}\StringTok{\textquotesingle{}Ps\textquotesingle{}}\NormalTok{], o[}\StringTok{\textquotesingle{}Es\textquotesingle{}}\NormalTok{], N)}
    \ControlFlowTok{elif}\NormalTok{ j }\OperatorTok{==} \DecValTok{7}\NormalTok{:}
\NormalTok{        n[}\StringTok{\textquotesingle{}d\textquotesingle{}}\NormalTok{] }\OperatorTok{=}\NormalTok{ (H }\OperatorTok{/} \DecValTok{6}\NormalTok{) }\OperatorTok{*}\NormalTok{ (}
\NormalTok{            o[}\StringTok{\textquotesingle{}k1\textquotesingle{}}\NormalTok{] }\OperatorTok{+} \DecValTok{2}\OperatorTok{*}\NormalTok{o[}\StringTok{\textquotesingle{}k2\textquotesingle{}}\NormalTok{] }\OperatorTok{+} \DecValTok{2}\OperatorTok{*}\NormalTok{o[}\StringTok{\textquotesingle{}k3\textquotesingle{}}\NormalTok{] }\OperatorTok{+}\NormalTok{ o[}\StringTok{\textquotesingle{}k4\textquotesingle{}}\NormalTok{]}
\NormalTok{        )}
    \ControlFlowTok{elif}\NormalTok{ j }\OperatorTok{==} \DecValTok{8}\NormalTok{:}
\NormalTok{        n[}\StringTok{\textquotesingle{}Um\textquotesingle{}}\NormalTok{] }\OperatorTok{=}\NormalTok{ U }\OperatorTok{*}\NormalTok{ PH2}
\NormalTok{        n[}\StringTok{\textquotesingle{}Pm\textquotesingle{}}\NormalTok{] }\OperatorTok{=}\NormalTok{ P }\OperatorTok{+} \FloatTok{0.5} \OperatorTok{*}\NormalTok{ o[}\StringTok{\textquotesingle{}d\textquotesingle{}}\NormalTok{][}\DecValTok{0}\NormalTok{]}
\NormalTok{        n[}\StringTok{\textquotesingle{}Em\textquotesingle{}}\NormalTok{] }\OperatorTok{=}\NormalTok{ E }\OperatorTok{+} \FloatTok{0.5} \OperatorTok{*}\NormalTok{ o[}\StringTok{\textquotesingle{}d\textquotesingle{}}\NormalTok{][}\DecValTok{1}\NormalTok{]}
    \ControlFlowTok{elif}\NormalTok{ j }\OperatorTok{==} \DecValTok{9}\NormalTok{:}
\NormalTok{        n[}\StringTok{\textquotesingle{}dphi\textquotesingle{}}\NormalTok{] }\OperatorTok{=}\NormalTok{ H }\OperatorTok{*}\NormalTok{ gN(o[}\StringTok{\textquotesingle{}Um\textquotesingle{}}\NormalTok{], o[}\StringTok{\textquotesingle{}Pm\textquotesingle{}}\NormalTok{], o[}\StringTok{\textquotesingle{}Em\textquotesingle{}}\NormalTok{])}
    \ControlFlowTok{elif}\NormalTok{ j }\OperatorTok{==} \DecValTok{10}\NormalTok{:}
\NormalTok{        n[}\StringTok{\textquotesingle{}U\textquotesingle{}}\NormalTok{] }\OperatorTok{=}\NormalTok{ U }\OperatorTok{*}\NormalTok{ PH}
\NormalTok{        n[}\StringTok{\textquotesingle{}P\textquotesingle{}}\NormalTok{] }\OperatorTok{=}\NormalTok{ P }\OperatorTok{+}\NormalTok{ o[}\StringTok{\textquotesingle{}d\textquotesingle{}}\NormalTok{][}\DecValTok{0}\NormalTok{]}
\NormalTok{        n[}\StringTok{\textquotesingle{}E\textquotesingle{}}\NormalTok{] }\OperatorTok{=}\NormalTok{ E }\OperatorTok{+}\NormalTok{ o[}\StringTok{\textquotesingle{}d\textquotesingle{}}\NormalTok{][}\DecValTok{1}\NormalTok{]}
\NormalTok{        n[}\StringTok{\textquotesingle{}H\textquotesingle{}}\NormalTok{] }\OperatorTok{=}\NormalTok{ o[}\StringTok{\textquotesingle{}H\textquotesingle{}}\NormalTok{] }\OperatorTok{*} \BuiltInTok{complex}\NormalTok{(}
\NormalTok{            math.cos(o[}\StringTok{\textquotesingle{}dphi\textquotesingle{}}\NormalTok{]), math.sin(o[}\StringTok{\textquotesingle{}dphi\textquotesingle{}}\NormalTok{])}
\NormalTok{        )}
\NormalTok{        n[}\StringTok{\textquotesingle{}Q\textquotesingle{}}\NormalTok{] }\OperatorTok{=}\NormalTok{ o[}\StringTok{\textquotesingle{}Q\textquotesingle{}}\NormalTok{] }\OperatorTok{*}\NormalTok{ math.exp(o[}\StringTok{\textquotesingle{}d\textquotesingle{}}\NormalTok{][}\DecValTok{2}\NormalTok{] }\OperatorTok{/}\NormalTok{ TAU0)}

    \CommentTok{\# sixth persistent state: same one{-}hot interaction, identity row}
\NormalTok{    Nc }\OperatorTok{=}\NormalTok{ [}\FloatTok{1.0} \OperatorTok{*}\NormalTok{ N }\ControlFlowTok{for}\NormalTok{ \_ }\KeywordTok{in} \BuiltInTok{range}\NormalTok{(}\DecValTok{11}\NormalTok{)]}
\NormalTok{    n[}\StringTok{\textquotesingle{}N\textquotesingle{}}\NormalTok{] }\OperatorTok{=} \BuiltInTok{float}\NormalTok{(blend(o[}\StringTok{\textquotesingle{}q\textquotesingle{}}\NormalTok{], Nc))}

\NormalTok{    n[}\StringTok{\textquotesingle{}q\textquotesingle{}}\NormalTok{] }\OperatorTok{=}\NormalTok{ np.roll(o[}\StringTok{\textquotesingle{}q\textquotesingle{}}\NormalTok{], }\DecValTok{1}\NormalTok{)}
    \ControlFlowTok{return}\NormalTok{ n}
\end{Highlighting}
\end{Shaded}

In the code,

\begin{Shaded}
\begin{Highlighting}[]
\NormalTok{j }\OperatorTok{=} \BuiltInTok{int}\NormalTok{(np.argmax(s[}\StringTok{\textquotesingle{}q\textquotesingle{}}\NormalTok{]))}
\end{Highlighting}
\end{Shaded}

is an implementation shortcut that selects the nonzero component of the
one-hot state quickly. The definition of the interaction in the formulas
is not an external loop counter but

\[
\sum_{j=0}^{10}q_jF_j.
\]

One macro step was verified with

\begin{Shaded}
\begin{Highlighting}[]
\KeywordTok{def}\NormalTok{ macro\_fast6(s):}
    \ControlFlowTok{for}\NormalTok{ \_ }\KeywordTok{in} \BuiltInTok{range}\NormalTok{(}\DecValTok{11}\NormalTok{):}
\NormalTok{        s }\OperatorTok{=}\NormalTok{ micro\_fast6(s)}
    \ControlFlowTok{return}\NormalTok{ s}
\end{Highlighting}
\end{Shaded}

Here too, the count of the \texttt{for} loop is not used as an input of
the state transition. Only \(q\) is used for the selection of the
action.

\begin{center}\rule{0.5\linewidth}{0.5pt}\end{center}

\subsection{7.3 The formulas of each internal
phase}\label{the-formulas-of-each-internal-phase}

Let \(h=\Delta\chi\) and

\[
\Phi_{1/2}=e^{ih/2},
\qquad
\Phi_1=e^{ih}.
\]

The first evaluation is

\[
K_1=R(U,P,E,\mathcal N),
\]

the second evaluation point is

\[
U_2=U\Phi_{1/2},
\]

\[
P_2=P+\frac h2 K_{1p},
\qquad
E_2=E+\frac h2 K_{1e},
\]

and the second evaluation is

\[
K_2=R(U_2,P_2,E_2,\mathcal N).
\]

Likewise

\[
U_3=U\Phi_{1/2},
\quad
P_3=P+\frac h2K_{2p},
\quad
E_3=E+\frac h2K_{2e},
\]

\[
K_3=R(U_3,P_3,E_3,\mathcal N),
\]

\[
U_4=U\Phi_1,
\quad
P_4=P+hK_{3p},
\quad
E_4=E+hK_{3e},
\]

\[
K_4=R(U_4,P_4,E_4,\mathcal N).
\]

The composition is

\[
\boxed{
D
=\frac h6
(K_1+2K_2+2K_3+K_4).
}
\]

Taking the angular evaluation point as

\[
U_m=U\Phi_{1/2},
\qquad
P_m=P+\frac12D_p,
\qquad
E_m=E+\frac12D_e,
\]

and

\[
\Delta\phi=h\,g_N(U_m,P_m,E_m),
\]

finally

\[
\boxed{
U'=U\Phi_1
}
\]

\[
\boxed{
P'=P+D_p,
\qquad
E'=E+D_e
}
\]

\[
\boxed{
H'=He^{i\Delta\phi}
}
\]

\[
\boxed{
Q'=Qe^{D_t/\tau_0}
}
\]

\[
\boxed{
\mathcal N'=\mathcal N
}
\]

This is not a sequential assignment such as \texttt{P\ +=\ Dp}. \(P'\)
and \(E'\) are sums of products written at once from the old state at
the start of the microstep alone.

\begin{center}\rule{0.5\linewidth}{0.5pt}\end{center}

\section{8. Full expansion of the interaction
kernels}\label{full-expansion-of-the-interaction-kernels}

\subsection{\texorpdfstring{8.1 Reading the trigonometric quantities
from
\(U\)}{8.1 Reading the trigonometric quantities from U}}\label{reading-the-trigonometric-quantities-from-u}

From

\[
U=e^{i\chi}
\]

we have

\[
C_\chi:=\cos\chi
=\frac{U+U^{-1}}2,
\]

\[
S_\chi:=\sin\chi
=\frac{U-U^{-1}}{2i}.
\]

Therefore \(\chi=\arg U\) need not be recovered inside the state update.

\begin{center}\rule{0.5\linewidth}{0.5pt}\end{center}

\subsection{8.2 The 2.5PN radiation-reaction
kernel}\label{the-2.5pn-radiation-reaction-kernel}

The same leading 2.5PN radiation reaction as in Paper 4 is used. For the
background of PN theory see {[}10{]}; for the lineage of the
radiation-reaction expressions see {[}11{]}.

\[
d=1+EC_\chi,
\qquad
r=\frac Pd,
\qquad
\sqrt P=P^{1/2},
\]

\[
\dot r=\frac{ES_\chi}{\sqrt P},
\qquad
v_t=\frac d{\sqrt P},
\]

\[
v^2=\dot r^2+v_t^2,
\qquad
r^{-1}=\frac dP.
\]

With the radiation-reaction coefficients

\[
C_{\rm RR}=\frac85\mathcal N r^{-3},
\]

\[
A
=\dot r
\left(
18v^2+\frac23r^{-1}-25\dot r^2
\right),
\]

\[
B
=-\left(
6v^2-2r^{-1}-15\dot r^2
\right),
\]

we have

\[
\dot E_{\rm orb}
=C_{\rm RR}(A\dot r+Bv^2),
\]

\[
\dot p=2C_{\rm RR}Bp,
\]

\[
E_{\rm orb}=\frac{e^2-1}{2p},
\]

\[
\dot e
=\frac{p\dot E_{\rm orb}+E_{\rm orb}\dot p}{e}.
\]

Also

\[
\frac{dt}{d\chi}
=\frac{r^2}{\sqrt p}.
\]

The core of the actual order-verification code is the following.

\begin{Shaded}
\begin{Highlighting}[]
\KeywordTok{def}\NormalTok{ rr\_numeric(U, P, E, N):}
\NormalTok{    Cc }\OperatorTok{=} \FloatTok{0.5} \OperatorTok{*}\NormalTok{ (U }\OperatorTok{+} \FloatTok{1.0} \OperatorTok{/}\NormalTok{ U)}
\NormalTok{    Ss }\OperatorTok{=}\NormalTok{ (U }\OperatorTok{{-}} \FloatTok{1.0} \OperatorTok{/}\NormalTok{ U) }\OperatorTok{/}\NormalTok{ (}\OtherTok{2j}\NormalTok{)}

\NormalTok{    den }\OperatorTok{=} \FloatTok{1.0} \OperatorTok{+}\NormalTok{ E }\OperatorTok{*}\NormalTok{ Cc}
\NormalTok{    r }\OperatorTok{=}\NormalTok{ P }\OperatorTok{/}\NormalTok{ den}
\NormalTok{    sqP }\OperatorTok{=}\NormalTok{ P}\OperatorTok{**}\FloatTok{0.5}

\NormalTok{    rdot }\OperatorTok{=}\NormalTok{ E }\OperatorTok{*}\NormalTok{ Ss }\OperatorTok{/}\NormalTok{ sqP}
\NormalTok{    vt }\OperatorTok{=}\NormalTok{ den }\OperatorTok{/}\NormalTok{ sqP}
\NormalTok{    v2 }\OperatorTok{=}\NormalTok{ rdot }\OperatorTok{*}\NormalTok{ rdot }\OperatorTok{+}\NormalTok{ vt }\OperatorTok{*}\NormalTok{ vt}
\NormalTok{    invr }\OperatorTok{=}\NormalTok{ den }\OperatorTok{/}\NormalTok{ P}

\NormalTok{    Crr }\OperatorTok{=} \FloatTok{1.6} \OperatorTok{*}\NormalTok{ N }\OperatorTok{*}\NormalTok{ invr}\OperatorTok{**}\DecValTok{3}

\NormalTok{    an }\OperatorTok{=}\NormalTok{ rdot }\OperatorTok{*}\NormalTok{ (}
        \FloatTok{18.0} \OperatorTok{*}\NormalTok{ v2}
        \OperatorTok{+}\NormalTok{ (}\FloatTok{2.0} \OperatorTok{/} \FloatTok{3.0}\NormalTok{) }\OperatorTok{*}\NormalTok{ invr}
        \OperatorTok{{-}} \FloatTok{25.0} \OperatorTok{*}\NormalTok{ rdot }\OperatorTok{*}\NormalTok{ rdot}
\NormalTok{    )}

\NormalTok{    bv }\OperatorTok{=} \OperatorTok{{-}}\NormalTok{(}
        \FloatTok{6.0} \OperatorTok{*}\NormalTok{ v2}
        \OperatorTok{{-}} \FloatTok{2.0} \OperatorTok{*}\NormalTok{ invr}
        \OperatorTok{{-}} \FloatTok{15.0} \OperatorTok{*}\NormalTok{ rdot }\OperatorTok{*}\NormalTok{ rdot}
\NormalTok{    )}

\NormalTok{    Edot }\OperatorTok{=}\NormalTok{ Crr }\OperatorTok{*}\NormalTok{ (an }\OperatorTok{*}\NormalTok{ rdot }\OperatorTok{+}\NormalTok{ bv }\OperatorTok{*}\NormalTok{ v2)}
\NormalTok{    pdot }\OperatorTok{=} \FloatTok{2.0} \OperatorTok{*}\NormalTok{ Crr }\OperatorTok{*}\NormalTok{ bv }\OperatorTok{*}\NormalTok{ P}

\NormalTok{    Eorb }\OperatorTok{=}\NormalTok{ (E }\OperatorTok{*}\NormalTok{ E }\OperatorTok{{-}} \FloatTok{1.0}\NormalTok{) }\OperatorTok{/}\NormalTok{ (}\FloatTok{2.0} \OperatorTok{*}\NormalTok{ P)}
\NormalTok{    edot }\OperatorTok{=}\NormalTok{ (P }\OperatorTok{*}\NormalTok{ Edot }\OperatorTok{+}\NormalTok{ Eorb }\OperatorTok{*}\NormalTok{ pdot) }\OperatorTok{/}\NormalTok{ E}

\NormalTok{    dt }\OperatorTok{=}\NormalTok{ r }\OperatorTok{*}\NormalTok{ r }\OperatorTok{/}\NormalTok{ sqP}

    \ControlFlowTok{return}\NormalTok{ pdot }\OperatorTok{*}\NormalTok{ dt, edot }\OperatorTok{*}\NormalTok{ dt, dt}
\end{Highlighting}
\end{Shaded}

This function takes the old-state \(U,P,E,\mathcal N\) as input and
returns \(dp/d\chi,de/d\chi,dt/d\chi\).

\begin{center}\rule{0.5\linewidth}{0.5pt}\end{center}

\subsection{8.3 The angular-progression
kernel}\label{the-angular-progression-kernel}

From the Darwin parameterization of bound Schwarzschild orbits {[}8{]},

\[
\frac{d\phi}{d\chi}
=
\left[
1-\frac{2(3+e\cos\chi)}p
\right]^{-1/2}
\]

is used.

Putting

\[
u
:=
\frac{3+EC_\chi}{P},
\]

we have

\[
(1-2u)^{-1/2}
=
\sum_{m=0}^{\infty}c_mu^m,
\]

\[
c_0=1,
\qquad
c_{m+1}=\frac{2m+1}{m+1}c_m.
\]

At order \(N\),

\[
\boxed{
g_N(U,P,E)
=\sum_{m=0}^{N}c_m
\left[
\frac{3+\frac E2(U+U^{-1})}{P}
\right]^m.
}
\]

In the code it is evaluated by Horner's method:

\begin{Shaded}
\begin{Highlighting}[]
\KeywordTok{def}\NormalTok{ gN\_numeric(U, P, E, N):}
\NormalTok{    C }\OperatorTok{=} \FloatTok{0.5} \OperatorTok{*}\NormalTok{ (U }\OperatorTok{+} \FloatTok{1.0} \OperatorTok{/}\NormalTok{ U)}
\NormalTok{    u }\OperatorTok{=}\NormalTok{ (}\FloatTok{3.0} \OperatorTok{+}\NormalTok{ E }\OperatorTok{*}\NormalTok{ C) }\OperatorTok{/}\NormalTok{ P}

\NormalTok{    c }\OperatorTok{=}\NormalTok{ COEFF[N]}
\NormalTok{    y }\OperatorTok{=} \BuiltInTok{complex}\NormalTok{(c[}\OperatorTok{{-}}\DecValTok{1}\NormalTok{])}
    \ControlFlowTok{for}\NormalTok{ a }\KeywordTok{in}\NormalTok{ c[}\OperatorTok{{-}}\DecValTok{2}\NormalTok{::}\OperatorTok{{-}}\DecValTok{1}\NormalTok{]:}
\NormalTok{        y }\OperatorTok{=} \BuiltInTok{complex}\NormalTok{(a) }\OperatorTok{+}\NormalTok{ u }\OperatorTok{*}\NormalTok{ y}
    \ControlFlowTok{return}\NormalTok{ y}
\end{Highlighting}
\end{Shaded}

\begin{center}\rule{0.5\linewidth}{0.5pt}\end{center}

\subsection{8.4 Comparison of the full expansion with the auxiliary
functions
removed}\label{comparison-of-the-full-expansion-with-the-auxiliary-functions-removed}

In an intermediate stage of the work, the stage that still kept
\texttt{rr()} and \texttt{gN()} was mistakenly called the ``full
expansion''. It was not a full expansion.

In the final audit, \(dp/d\chi\) and \(de/d\chi\) were expanded into
finite Laurent polynomials in \(U^{\pm1}\), \(g_{12}\) was also written
out with explicit coefficients, and they were compared with the original
formulas at 100,000 random points.

The maximum errors were

\[
|\Delta(dp/d\chi)|_{\max}
\approx1.25\times10^{-16},
\]

\[
|\Delta(de/d\chi)|_{\max}
\approx2.08\times10^{-17},
\]

\[
\left|\frac{\Delta(dt/d\chi)}{dt/d\chi}\right|_{\max}
\approx8.61\times10^{-16},
\]

\[
|\Delta g_{12}|_{\max}
\approx8.88\times10^{-16},
\]

in agreement within floating-point rounding. {[}Numerical{]}

\begin{center}\rule{0.5\linewidth}{0.5pt}\end{center}

\section{9. Readout}\label{readout}

\subsection{9.1 Orbital position}\label{orbital-position}

From the states \(U,P,E,H\), compute

\[
C_\chi=\frac{U+U^{-1}}2
\]

and take

\[
\boxed{
r=\frac{P}{1+EC_\chi}.
}
\]

Since

\[
H=e^{i\phi},
\]

we have

\[
\boxed{
x=r\operatorname{Re}H,
\qquad
y=r\operatorname{Im}H.
}
\]

\begin{center}\rule{0.5\linewidth}{0.5pt}\end{center}

\subsection{9.2 Clock}\label{clock}

Since

\[
Q=e^{t/\tau_0},
\]

only at observation is it read as

\[
\boxed{
t_{\rm obs}
=\tau_0\log |Q|.
}
\]

The \(t_{\rm obs}\) obtained here is not returned to the state update.

Organized as a function, the readout operation in the audit code is the
following. It is equivalent to the operation actually used in the audit.

\begin{Shaded}
\begin{Highlighting}[]
\KeywordTok{def}\NormalTok{ readout(U, P, E, H, Q, tau0):}
\NormalTok{    cos\_chi }\OperatorTok{=} \FloatTok{0.5} \OperatorTok{*}\NormalTok{ (U }\OperatorTok{+} \FloatTok{1.0} \OperatorTok{/}\NormalTok{ U)}
\NormalTok{    r }\OperatorTok{=}\NormalTok{ (P }\OperatorTok{/}\NormalTok{ (}\FloatTok{1.0} \OperatorTok{+}\NormalTok{ E }\OperatorTok{*}\NormalTok{ cos\_chi)).real}

\NormalTok{    x }\OperatorTok{=}\NormalTok{ r }\OperatorTok{*}\NormalTok{ H.real}
\NormalTok{    y }\OperatorTok{=}\NormalTok{ r }\OperatorTok{*}\NormalTok{ H.imag}

    \CommentTok{\# log is used only here, never in the state update}
\NormalTok{    t }\OperatorTok{=}\NormalTok{ tau0 }\OperatorTok{*}\NormalTok{ math.log(}\BuiltInTok{abs}\NormalTok{(Q))}

    \ControlFlowTok{return}\NormalTok{ x, y, r, t}
\end{Highlighting}
\end{Shaded}

The fundamental flow of information of this paper is therefore separated
into

\[
\boxed{
(U,P,E,H,Q)
\xrightarrow{\text{interaction}}
(U',P',E',H',Q')
}
\]

and

\[
\boxed{
(U,P,E,H,Q)
\xrightarrow{\text{readout}}
(x,y,r,t).
}
\]

\begin{center}\rule{0.5\linewidth}{0.5pt}\end{center}

\section{10. Program audit results for scheme B and the sixth
state}\label{program-audit-results-for-scheme-b-and-the-sixth-state}

For scheme B, the rules as a state machine were audited directly, not
merely whether the final orbit looks similar.

For the 11-phase machine of the first audit, it was checked

\begin{enumerate}
\def\labelenumi{\arabic{enumi}.}
\tightlist
\item
  whether there is any place where the updated state is read within the
  same microstep,
\item
  whether the one-hot phase is ever broken,
\item
  whether the persistent state changes illegitimately in any phase other
  than the commit phase,
\item
  whether the phase returns to the original after 11 microsteps,
\item
  whether one macro step coincides with the original Paper 4 RK4 step.
\end{enumerate}

In the one-macro-step comparison with 500 random states, for the five
components of the first audit,

\[
\boxed{
\max|\Delta(U,P,E,H,Q)|=(0,0,0,0,0).
}
\]

In the long comparison over 12 conditions,

\[
\boxed{
\max|\Delta(U,P,E,H,Q)|
=
(0,0,0,3.38\times10^{-16},4.00\times10^{-15}).
}
\]

In the additional \(\nu\) audit, over 12 conditions \(\times\) 4000
macrosteps the first five components of the old 5-state and the new
6-state systems, and \(r,x,y,t\), coincided down to the bit patterns of
the stored values, and the bitwise mismatch of \(\mathcal N\) was also
0. Over 3 representative conditions \(\times\) 12000 macrosteps the
maximum difference of all compared quantities was 0, and there were 0
violations of the microstep identity holding of \(\mathcal N\).

The maximum absolute differences against the stored C1 CSV are identical
for the old and the new systems:

\[
|\Delta p|_{\max}=4.97\times10^{-14},
\qquad
|\Delta e|_{\max}=7.77\times10^{-16},
\]

\[
|\Delta t|_{\max}=1.94\times10^{-10},
\]

\[
|\Delta x|_{\max}=3.85\times10^{-12},
\qquad
|\Delta y|_{\max}=5.11\times10^{-12}.
\]

These tiny differences did not arise from the present internalization of
\(\nu\); they are the differences in rounding and implementation path
that the old scheme-B audit system already had against the stored C1.
{[}Numerical{]}

\begin{figure}
\centering
\pandocbounded{\includegraphics[keepaspectratio,alt={Figure 4. Strict audit of scheme B over 12 conditions}]{figs/figure03.png}}
\caption{Figure 4. Strict audit of scheme B over 12 conditions}
\end{figure}

\textbf{Figure 4.} The difference between scheme B and the reference
update over 12 conditions. Even the nonzero components are essentially
within floating-point rounding. In the control of the internalization of
\(\nu\), the first five components of the old 5-state and the new
6-state systems are bitwise identical, and the sixth state
\(\mathcal N\) is also held as an identity.

From the above, for the present limited model it was confirmed that

\begin{center}
\fbox{\parbox{0.90\linewidth}{\centering the sequential processing of Paper 4 can be rearranged into an explicit state machine without the external $\nu.$}}
\end{center}

This does not mean, however, that the 11-phase state machine is a
fundamental action of nature, or that the whole existing scheme B
satisfies the strictest fully multiplicative internal operation
condition defined afterwards.

\section{11. Separating the approximation
orders}\label{separating-the-approximation-orders}

Several ``orders'' exist in Paper 4 and in this paper. They must not be
confused.

\subsubsection{Table 2. Distinction between the approximation orders and
the state
closure}\label{table-2.-distinction-between-the-approximation-orders-and-the-state-closure}

{\def\LTcaptype{none} % do not increment counter
\begin{longtable}[]{@{}
  >{\raggedright\arraybackslash}p{(\linewidth - 8\tabcolsep) * \real{0.2000}}
  >{\raggedright\arraybackslash}p{(\linewidth - 8\tabcolsep) * \real{0.2000}}
  >{\raggedright\arraybackslash}p{(\linewidth - 8\tabcolsep) * \real{0.2000}}
  >{\raggedright\arraybackslash}p{(\linewidth - 8\tabcolsep) * \real{0.2000}}
  >{\raggedright\arraybackslash}p{(\linewidth - 8\tabcolsep) * \real{0.2000}}@{}}
\toprule\noalign{}
\begin{minipage}[b]{\linewidth}\raggedright
Layer
\end{minipage} & \begin{minipage}[b]{\linewidth}\raggedright
Present / verified order
\end{minipage} & \begin{minipage}[b]{\linewidth}\raggedright
Error or unretained terms
\end{minipage} & \begin{minipage}[b]{\linewidth}\raggedright
Meaning
\end{minipage} & \begin{minipage}[b]{\linewidth}\raggedright
Relation to the number of states
\end{minipage} \\
\midrule\noalign{}
\endhead
\bottomrule\noalign{}
\endlastfoot
Dissipative physical model & leading 2.5PN radiation reaction & Higher
PN terms not introduced & Approximation order of the physical model & A
separate issue from the six persistent states \\
Angular series & \(N=12\) baseline; tested \(12,16,20,24\) &
\(O(u^{N+1})\) & Series truncation & Raised without increasing the
number of states \\
Paper 4 time integration & RK4 & local \(O(h^5)\), global \(O(h^4)\) &
Order of numerical integration & The 11-phase expansion makes the
processing order explicit \\
B 11-phase expansion & Algebraic rearrangement of RK4 & No new
truncation error & Machine-state depth & Not a claim of a physical
11-stage structure \\
Taylor order verification & 4, 6, 8 & Measured convergence confirmed
separately & Depth of numerical approximation & Raised without
increasing the number of persistent states \\
Internalization of \(\nu\) & \(\mathcal N=\nu\) identity state & Not an
approximation & Correction of the state closure & 5 to 6 is not a
precision improvement but the removal of an information leakage \\
\end{longtable}
}

In particular,

\begin{center}
\fbox{\parbox{0.90\linewidth}{\centering raising the order of the numerical integration$\neq$raising the PN physical order$.$}}
\end{center}

Whether

\[
\boxed{
\text{raising the approximation order}
\neq
\text{adding a physical state}
}
\]

holds is verified next.

\begin{center}\rule{0.5\linewidth}{0.5pt}\end{center}

\section{12. Raising the order of the angular
series}\label{raising-the-order-of-the-angular-series}

The exact expression of the angular factor is

\[
g_\infty(u)=(1-2u)^{-1/2},
\]

and at finite \(N\)

\[
g_N(u)=\sum_{m=0}^{N}c_mu^m.
\]

With the persistent states \(Z_6=(U,P,E,H,Q,\mathcal N)\) fixed, only
the truncation order of the angular series was increased to

\[
N=12,16,20,24.
\]

For one of the most severe test conditions,

\[
p_0=24.3,
\qquad e_0=0.55,
\]

the azimuth error after three orbits decreased as

{\def\LTcaptype{none} % do not increment counter
\begin{longtable}[]{@{}rr@{}}
\toprule\noalign{}
\(N\) & \(|\Delta\phi|\) \\
\midrule\noalign{}
\endhead
\bottomrule\noalign{}
\endlastfoot
12 & \(1.56\times10^{-7}\) \\
16 & \(9.03\times10^{-10}\) \\
20 & \(5.54\times10^{-12}\) \\
24 & \(3.55\times10^{-14}\) \\
\end{longtable}
}

\begin{figure}
\centering
\pandocbounded{\includegraphics[keepaspectratio,alt={Figure 5. Order of the angular series and the error}]{figs/figure04.png}}
\caption{Figure 5. Order of the angular series and the error}
\end{figure}

\textbf{Figure 5.} Convergence for \(N=12,16,20,24\) over 9 conditions.
The number of states is not changed; only the truncation position of the
same series generating rule is changed.

From this result, at least in a finite range,

\[
\boxed{
N\uparrow
\quad\text{while}\quad
\dim Z_6=6
}
\]

holds.

\begin{center}\rule{0.5\linewidth}{0.5pt}\end{center}

\section{13. Can the order of the numerical integration be raised with
the six persistent states
kept?}\label{can-the-order-of-the-numerical-integration-be-raised-with-the-six-persistent-states-kept}

\subsection{13.1 Direct Taylor
recurrence}\label{direct-taylor-recurrence}

Increasing the number of RK4 stages increases the internal stage states
as well.

In order to examine separately ``whether a new persistent physical state
is needed to raise the order'', the Taylor coefficients were generated
within a single step for the five evolving components

\[
Z=(U,P,E,H,Q)
\]

of the same six persistent states, with \(\mathcal N\) held as an
identity.

Writing the continuous vector field as

\[
\frac{dZ}{d\chi}=F_N(Z),
\]

we take

\[
Z(\chi+h)
=\sum_{k=0}^{p}A_k h^k+O(h^{p+1}),
\]

\[
A_0=Z_n,
\]

\[
\boxed{
(k+1)A_{k+1}
=[\xi^k]
F_N\!\left(\sum_{j=0}^{k}A_j\xi^j\right).
}
\]

The core of the code is the following.

\begin{Shaded}
\begin{Highlighting}[]
\KeywordTok{def}\NormalTok{ taylor\_step(z, h, order, N, tau0):}
    \CommentTok{\# Same five dynamically evolving components.}
    \CommentTok{\# Revised closure appends N as an identity persistent state.}
    \CommentTok{\# No RK stage is promoted to a new persistent state.}
\NormalTok{    a }\OperatorTok{=}\NormalTok{ np.zeros((}\DecValTok{5}\NormalTok{, order }\OperatorTok{+} \DecValTok{1}\NormalTok{), }\BuiltInTok{complex}\NormalTok{)}
\NormalTok{    a[:, }\DecValTok{0}\NormalTok{] }\OperatorTok{=}\NormalTok{ z}

    \ControlFlowTok{for}\NormalTok{ k }\KeywordTok{in} \BuiltInTok{range}\NormalTok{(order):}
\NormalTok{        series }\OperatorTok{=}\NormalTok{ [TPS(a[j], order) }\ControlFlowTok{for}\NormalTok{ j }\KeywordTok{in} \BuiltInTok{range}\NormalTok{(}\DecValTok{5}\NormalTok{)]}
\NormalTok{        U, P, E, H, Q }\OperatorTok{=}\NormalTok{ series}

\NormalTok{        dp, de, dt }\OperatorTok{=}\NormalTok{ rr\_tps(U, P, E)}
\NormalTok{        g }\OperatorTok{=}\NormalTok{ gN\_tps(U, P, E, N)}

\NormalTok{        F }\OperatorTok{=}\NormalTok{ [}
            \OtherTok{1j} \OperatorTok{*}\NormalTok{ U,}
\NormalTok{            dp,}
\NormalTok{            de,}
            \OtherTok{1j} \OperatorTok{*}\NormalTok{ g }\OperatorTok{*}\NormalTok{ H,}
\NormalTok{            (dt }\OperatorTok{/}\NormalTok{ tau0) }\OperatorTok{*}\NormalTok{ Q,}
\NormalTok{        ]}

        \ControlFlowTok{for}\NormalTok{ j }\KeywordTok{in} \BuiltInTok{range}\NormalTok{(}\DecValTok{5}\NormalTok{):}
\NormalTok{            a[j, k }\OperatorTok{+} \DecValTok{1}\NormalTok{] }\OperatorTok{=}\NormalTok{ F[j].c[k] }\OperatorTok{/}\NormalTok{ (k }\OperatorTok{+} \DecValTok{1}\NormalTok{)}

\NormalTok{    powers }\OperatorTok{=}\NormalTok{ np.array(}
\NormalTok{        [h}\OperatorTok{**}\NormalTok{k }\ControlFlowTok{for}\NormalTok{ k }\KeywordTok{in} \BuiltInTok{range}\NormalTok{(order }\OperatorTok{+} \DecValTok{1}\NormalTok{)],}
        \BuiltInTok{complex}
\NormalTok{    )}
    \ControlFlowTok{return}\NormalTok{ a }\OperatorTok{@}\NormalTok{ powers}
\end{Highlighting}
\end{Shaded}

Here the Taylor coefficients \(A_k\) are computational coefficients
within one step, and are not carried to the next macro step.

In this numerical convergence experiment only the five evolving
components are arrayed, but in the revised closure \(\mathcal N\) is
placed alongside as an identity state. Therefore \(p=4,6,8\) can be
compared with the six persistent states kept and the physical formulas
unchanged.

\begin{center}\rule{0.5\linewidth}{0.5pt}\end{center}

\subsection{13.2 Measured convergence
orders}\label{measured-convergence-orders}

Over 9 conditions the steps per orbit were varied as 8, 16, 32 and 64.

Looking at the convergence rate from 16 to 32 steps per orbit, where the
influence of rounding error is relatively small, averaged over the 9
conditions:

\subsubsection{Table 3. Average measured order from 16 to 32 steps per
orbit}\label{table-3.-average-measured-order-from-16-to-32-steps-per-orbit}

{\def\LTcaptype{none} % do not increment counter
\begin{longtable}[]{@{}rrrr@{}}
\toprule\noalign{}
Design order & \(P\) & \(E\) & \(\phi\) \\
\midrule\noalign{}
\endhead
\bottomrule\noalign{}
\endlastfoot
4 & 3.998 & 3.993 & 3.928 \\
6 & 5.997 & 6.010 & 5.921 \\
8 & 7.918 & 8.072 & 7.916 \\
\end{longtable}
}

At orders 4 and 6 the design order is almost kept from 32 to 64 as well.
At order 8, \(P,E\) begin to approach machine precision, so that at
finer steps the influence of the reference-solution error and of
rounding error can no longer be neglected.

\begin{figure}
\centering
\pandocbounded{\includegraphics[keepaspectratio,alt={Figure 6. Convergence at orders 4, 6 and 8 with the same persistent states}]{figs/figure05.png}}
\caption{Figure 6. Convergence at orders 4, 6 and 8 with the same
persistent states}
\end{figure}

\textbf{Figure 6.} Convergence test in which only the approximation
depth is changed for the same vector field and the same five evolving
components. The sixth state \(\mathcal N\) after the revision is held as
an identity and is independent of this raising of the order.

The overall picture over the 9 conditions is shown in Figure 7.

\begin{figure}
\centering
\pandocbounded{\includegraphics[keepaspectratio,alt={Figure 7. Comparison of the raising of the order over 9 conditions}]{figs/figure06.png}}
\caption{Figure 7. Comparison of the raising of the order over 9
conditions}
\end{figure}

\textbf{Figure 7.} The decrease of the error with the raising of the
order is confirmed in strong-, intermediate- and weak-field conditions
and at different eccentricities.

What was confirmed in this paper is therefore

\begin{center}
\fbox{\parbox{0.90\linewidth}{\centering increase of the approximation order$\not\Rightarrow$addition of a new persistent physical state$.$}}
\end{center}

The amount of computation and the number of temporary Taylor
coefficients, however, increase with the order.

\begin{center}\rule{0.5\linewidth}{0.5pt}\end{center}

\section{\texorpdfstring{14. On the fixed action \(S\) and the sequence
of state-dependent
actions}{14. On the fixed action S and the sequence of state-dependent actions}}\label{on-the-fixed-action-s-and-the-sequence-of-state-dependent-actions}

In the audit of this paper, when the present representation of Paper 4
was expanded as it is into explicit states, the processing could not be
expressed by a single fixed action, and several internal actions
depending on the processing phase were used.

In the formulas,

\[
\boxed{
\mathcal Z_{m+1}
=F_{\sigma_m}(\mathcal Z_m),
\qquad
\sigma_m\ \text{is encoded in }q_m.
}
\]

From this result, however,

\[
\boxed{
S_n\text{ is always necessary in nature}
}
\]

cannot be concluded.

There are three reasons.

First, choosing other state coordinates may make it possible to unify
into a single fixed action.

Second, the 11 internal phases make explicit the processing structure of
the numerical algorithm RK4, and do not prove an 11-stage structure of
the physics itself.

Third, in the Taylor recurrence the order could be raised with the same
set of persistent states and a single vector field \(F_N\).

The conclusion of this paper is therefore

\begin{center}
\fbox{\parbox{0.90\linewidth}{\centering a sequence of actions was necessary in the present B construction, but its necessity in principle is unproven$.$}}
\end{center}

\begin{center}\rule{0.5\linewidth}{0.5pt}\end{center}

\section{15. What was learned from this
paper}\label{what-was-learned-from-this-paper}

\subsection{15.1 Self-correction of the ``two-state
model''}\label{self-correction-of-the-two-state-model}

The \(X=(a,b)\) of Paper 4 was a useful readout representation, but it
was not a fundamental state containing all the information that
determines the next state. This does not invalidate the numerical
results of Paper 4. Paper 5 makes explicit which upstream information
those results depended on.

\subsection{15.2 The first audit's ``closed with five states'' was
corrected
again}\label{the-first-audits-closed-with-five-states-was-corrected-again}

The first audit judged that the system closed once \(Z_5=(U,P,E,H,Q)\)
was made explicit. But tracking the radiation-reaction kernel showed
that \(\nu\) was supplied from outside.

This paper itself therefore corrected the intermediate conclusion of the
audit to

\[
\boxed{
Z_5'=F(Z_5;\nu)
\quad\longrightarrow\quad
Z_6'=\widetilde F(Z_6).
}
\]

What matters here is not the number ``6'' but the audit procedure:
\textbf{the information that determines the next state is not hidden as
a computational auxiliary value, and when found, it is moved to a state
candidate}.

\subsection{15.3 Making the hidden intermediate update values explicit
yields a state
machine}\label{making-the-hidden-intermediate-update-values-explicit-yields-a-state-machine}

In ordinary sequential numerical code, the evaluation points and the
intermediate values of the integration are treated as ``matters of
computational convenience''. Under the rules of this paper, however, as
long as information is held for use in the next operation, it is treated
as an explicit work/machine state. As a result, the update of Paper 4
could be rewritten as an 11-phase state machine with a one-hot phase.

\subsection{15.4 The readout and the state transition can be
separated}\label{the-readout-and-the-state-transition-can-be-separated}

For the clock in particular, taking

\[
Q=e^{t/\tau_0}
\]

as the state, neither \(t\) nor \(\log Q\) is needed during the state
update.

\[
t=\tau_0\log|Q|
\]

is used only at readout. Likewise \(x,y,r\) are readouts and are not
returned to the state update.

\begin{center}
\fbox{\parbox{0.90\linewidth}{\centering the operational form of the internal state$\neq$the coordinate form read by the observer}}
\end{center}

could be implemented concretely.

\subsection{15.5 Raising the precision and correcting the state closure
can be
distinguished}\label{raising-the-precision-and-correcting-the-state-closure-can-be-distinguished}

Raising the order of the angular series and of the numerical integration
did not require adding a persistent state. On the other hand, the
internalization of \(\nu\) increased the number of states from 5 to 6 in
order to remove an external information leakage, independently of any
precision improvement.

Therefore, at least in this model,

\[
\boxed{
\text{raising the precision}
\neq
\text{correcting the state closure}.
}
\]

\section{16. Limits}\label{limits}

This paper has the following limits.

\begin{enumerate}
\def\labelenumi{\arabic{enumi}.}
\item
  \textbf{The minimality of the six persistent states is unproven.}\\
  Using other coordinates, conserved quantities, or a more fundamental
  state representation, the system might close with fewer persistent
  states.
\item
  \textbf{The minimality of the 11-phase machine state and its necessity
  in nature are unproven.}\\
  The 11 phases are the structure of the present implementation that
  expands RK4 without hidden intermediate values; it is not a claim that
  nature acts in 11 stages.
\item
  \textbf{The conformity of the whole existing scheme B to the ``fully
  multiplicative internal operation'' is unproven.}\\
  What was additionally confirmed in this paper is only that the old
  output is kept exactly when the external \(\nu\) is internalized as
  the sixth state \(\mathcal N\).
\item
  \textbf{The physical approximations of Paper 4 are inherited.}\\
  The conservative part is of the Schwarzschild / test-particle type and
  the dissipation is the leading 2.5PN {[}4,8-11{]}; this is not full
  finite-mass-ratio general relativity.
\item
  \textbf{Physical orders higher than 2.5PN are not verified.}\\
  Raising the numerical integration to order 8 does not raise the
  physical model itself beyond 2.5PN.
\item
  \textbf{The complex representation of \(U,H\) is an implementation
  choice.}\\
  Complex numbers are not claimed to be fundamental reality.
\item
  \textbf{Convergence of the whole physical model to the exact solution
  is not proved.}\\
  The angular series showed clear convergence to a known analytic
  function, but whether the infinite-order limit of the whole model
  coincides with exact GR is a separate problem.
\end{enumerate}

The point reached by this paper is therefore limited not to the
discovery of the minimal registers of nature or of a universal
interaction, but to \textbf{having carried out the state-closure audit
of the present two-body gravity model from 2 to 5 to 6 and having
removed the external \(\nu\)}.

\section{17. Reproducibility and research
record}\label{reproducibility-and-research-record}

Since in this paper the program audit itself is the experiment, not only
the final code but also the intermediate path was saved.

In the research folder on Google Drive (the folder names in the
repository are in Japanese; transliterated here),

\begin{Shaded}
\begin{Highlighting}[]
\NormalTok{anonymous\_vertex\_state\_generation\_geometry/}
\NormalTok{fifth\_thought\_experiment\_multiplicative\_state\_expansion\_and\_order\_verification\_20260924/}
\end{Highlighting}
\end{Shaded}

at least the following are saved.

\begin{itemize}
\tightlist
\item
  \texttt{00\_\allowbreak{}reference\_\allowbreak{}from\_\allowbreak{}paper4/\allowbreak{}} - reference code of Paper 4 and
  the C1 raw data
\item
  \texttt{01\_\allowbreak{}reconstruction\_\allowbreak{}of\_\allowbreak{}five\_\allowbreak{}state\_\allowbreak{}synchronous\_\allowbreak{}map/\allowbreak{}} -
  first audit: five persistent states
\item
  \texttt{02\_\allowbreak{}multiplicative\_\allowbreak{}state\_\allowbreak{}representation\_\allowbreak{}and\_\allowbreak{}synchronous\_\allowbreak{}inlining\_\allowbreak{}history/\allowbreak{}}
  - including intermediate and rejected proposals
\item
  \texttt{03\_\allowbreak{}A\_\allowbreak{}local\_\allowbreak{}exp\_\allowbreak{}log\_\allowbreak{}vs\_\allowbreak{}B\_\allowbreak{}internal\_\allowbreak{}state\_\allowbreak{}expansion/\allowbreak{}} -
  comparison of A and B
\item
  \texttt{04\_\allowbreak{}scheme\_\allowbreak{}B\_\allowbreak{}11\_\allowbreak{}phase\_\allowbreak{}synchronous\_\allowbreak{}state\_\allowbreak{}machine\_\allowbreak{}strict\_\allowbreak{}audit/\allowbreak{}}
  - static and numerical audit of scheme B
\item
  \texttt{05\_\allowbreak{}full\_\allowbreak{}expansion\_\allowbreak{}of\_\allowbreak{}operation\_\allowbreak{}blocks\_\allowbreak{}Laurent\_\allowbreak{}N12/\allowbreak{}} -
  Laurent expansion and the 100,000-point comparison
\item
  \texttt{06\_\allowbreak{}approximation\_\allowbreak{}order\_\allowbreak{}and\_\allowbreak{}state\_\allowbreak{}S\_\allowbreak{}non\_\allowbreak{}increase\_\allowbreak{}verification/\allowbreak{}}
  - \(N=12,16,20,24\) and orders 4, 6, 8
\item
  \texttt{08\_\allowbreak{}nu\_\allowbreak{}sixth\_\allowbreak{}state\_\allowbreak{}internalization\_\allowbreak{}control\_\allowbreak{}experiment\_\allowbreak{}20260925/\allowbreak{}}
  - internalization of the external \(\nu\) as the sixth state,
  12-condition and long-run controls, bitwise audit
\item
  \texttt{paper5\_\allowbreak{}revision\_\allowbreak{}v2\_\allowbreak{}figures\_\allowbreak{}tables\_\allowbreak{}20260926/\allowbreak{}} - revised
  figures and tables, source data, figure/table regeneration program,
  reproducibility ZIP
\end{itemize}

The revised figures and tables can be regenerated from the stored CSV /
JSON. The figure/table regeneration program is saved as

\begin{Shaded}
\begin{Highlighting}[]
\NormalTok{paper5\_revision\_v2\_figures\_tables\_20260926/scripts/}
\NormalTok{generate\_paper5\_revision\_v2.py}
\end{Highlighting}
\end{Shaded}

The main reproduction code for the internalization of \(\nu\) is

\begin{Shaded}
\begin{Highlighting}[]
\NormalTok{08\_nu\_sixth\_state\_internalization\_control\_experiment\_20260925/}
\NormalTok{run\_nu\_internal\_state\_control\_v1.py}
\end{Highlighting}
\end{Shaded}

The orbit figure of Figure 2 can be regenerated with

\begin{Shaded}
\begin{Highlighting}[]
\NormalTok{plot\_c1\_orbit\_6state\_audit\_v1.py}
\end{Highlighting}
\end{Shaded}

in the same folder. This program reads the stored C1 CSV for reference
comparison, separately generates 12000 macrosteps of the sixth-state
generator, and then overlays them. The reference orbit is not fed back
into the state update during the generation.

In the revision of this paper, the results of the charged eight-state
experiments (09 onwards) are not used. The fifth thought experiment
closes with the state-closure audit of the gravitational system alone.

\section{18. Conclusion}\label{conclusion}

The fourth thought experiment {[}4{]} was, on the surface, an experiment
reconstructing relativistic gravitational orbits from the interaction of
the two-component state

\[
X_{n+1}=S_nX_n.
\]

The first audit of this paper, however, made clear that the actual orbit
generation used the five persistent states

\[
\chi,p,e,\phi,t,
\]

and that \(X=(a,b)\) was a derived representation made from them. They
were therefore re-encoded as

\[
Z_5=(U,P,E,H,Q),
\]

and by expanding the intermediate values and the processing phases of
RK4 as explicit states as well, the sequential processing of Paper 4 was
reconstructed into an 11-phase synchronous state machine.

Stopping the audit there, however, was insufficient. Since the symmetric
mass ratio \(\nu\) remained as an external argument in the
radiation-reaction kernel, the system was in fact

\[
Z_5'=F(Z_5;\nu).
\]

Accordingly,

\[
\boxed{
Z_6=(U,P,E,H,Q,\mathcal N),
\qquad
\mathcal N=\nu
}
\]

was taken, and \(\mathcal N\) was internalized as an identity row of the
11-phase one-hot interaction. The new transition function has no
external \texttt{nu} argument.

Over 12 conditions \(\times\) 4000 macrosteps, the first five components
of the old 5-state and the new 6-state systems, and \(r,x,y,t\), were
bitwise identical, and the bitwise mismatch of \(\mathcal N\) was also
0. Over 3 representative conditions \(\times\) 12000 macrosteps the
maximum difference of the compared quantities was 0, and there were 0
violations of the microstep identity holding of \(\mathcal N\). The
present correction from 5 to 6 is therefore not a modification that
changes the physical formulas or the orbit, but \textbf{a correction of
the state closure that moves the place of physical information left
outside into the state}.

Furthermore, raising the truncation order of the angular series or the
order of the numerical integration did not require adding a new
persistent physical state. In this paper, therefore,

\[
\boxed{
\text{correction of the state closure}
\neq
\text{numerical precision improvement}
}
\]

could be separated concretely.

Written in the most restricted form, the point reached by this paper is
to have traced, with formulas, code and bitwise control experiments, the
path of self-correction

\[
\boxed{
\begin{array}{c}
\text{the apparent two states of Paper 4}\
\downarrow\\
\text{first audit: five persistent states}\
\downarrow\\
\text{additional audit: the external }\nu\text{ is found}\
\downarrow\\
\text{internalization into the six persistent states }Z_6
\end{array}
}
\]

This does not mean that six states are the minimal registers of nature,
that the 11 phases are the real structure of the fundamental
interaction, or that general relativity has been derived from six
states. What is established at present is only that, for this limited
system reverse-constructed from the known two-body gravity model,
\textbf{the information that determines the next state could be
rearranged, without leaving it as a hidden external parameter, into a
synchronous state map with six persistent logical states and explicit
machine/work states}.

\section{References}\label{references}

\subsection{This series}\label{this-series}

{[}0{]} Noriaki Kihara, \textbf{Designing a Research Project to Explore
the Emergence of Spacetime, Particles, and Interactions from States and
Relations: A Framework of Preliminary Design and Feasibility
Verification for Foundational-Physics Model Exploration under
Uncertainty}, v1.0 (2026-09-20). Concept DOI: 10.5281/zenodo.22851944;
Version DOI: 10.5281/zenodo.22851945.

{[}1{]} Noriaki Kihara, \textbf{The First Thought Experiment: Tracing
the Equivalence Principle Down to Two Nameless Degrees of Freedom -
Deriving the Sign of the Sum of Squares, the Curvature Radius, and Three
Systems (Inertial, Uniformly Accelerated, Rotating) from the
Conservation of an Area Readout}, v1.0 (2026-09-20). Concept DOI:
10.5281/zenodo.22857952; Version DOI: 10.5281/zenodo.22857953.

{[}2{]} Noriaki Kihara, \textbf{The Second Thought Experiment: Deriving
the Coulomb-Type Inverse-Square Force from a Product Readout and an Area
Clock - Keeping Two Values and a Linear Law, Obtaining the Focus,
Attraction and Repulsion and Neutrality, and Fixing What Two Values
Cannot Write}, v1.2 (2026-09-21). Concept DOI: 10.5281/zenodo.22867336;
Version DOI: 10.5281/zenodo.22876596.

{[}3{]} Noriaki Kihara, \textbf{The Third Thought Experiment: Reading
Out the State of Two Bodies a, b from the Orbit of Two Celestial Bodies
a, b - Obtaining the Mass Ratio, the Relative Distance and the Relative
Time from Orbit and Gravitational-Wave Readouts, with Two Values of
Undetermined Meaning and the Knowledge \(G=c=M=1\) Alone}, v1.5
(2026-09-23). Concept DOI: 10.5281/zenodo.22909660; Version DOI:
10.5281/zenodo.22909661.

{[}4{]} Noriaki Kihara, \textbf{The Fourth Thought Experiment: From a
Fixed Action \(S\) to a State-Dependent Action \(S_n\) - Can the
Relativistic Gravitational Orbits Computed Externally in the Third
Thought Experiment Be Regenerated from the Sequential Interaction
\(X_{n+1}=S_nX_n\)?}, v1.0 (2026-09-23). Concept DOI:
10.5281/zenodo.22919787; Version DOI: 10.5281/zenodo.22919788.

\subsection{External}\label{external}

{[}5{]} R. Zwanzig, ``Memory Effects in Irreversible Thermodynamics,''
\emph{Physical Review} \textbf{124}, 983-992 (1961). DOI:
10.1103/PhysRev.124.983.

{[}6{]} H. F. Trotter, ``On the Product of Semi-Groups of Operators,''
\emph{Proceedings of the American Mathematical Society} \textbf{10}(4),
545-551 (1959). DOI: 10.1090/S0002-9939-1959-0108732-6.

{[}7{]} J. C. Butcher, ``On Runge-Kutta Processes of High Order,''
\emph{Journal of the Australian Mathematical Society} \textbf{4}(2),
179-194 (1964). DOI: 10.1017/S1446788700023387.

{[}8{]} C. G. Darwin, ``The Gravity Field of a Particle,''
\emph{Proceedings of the Royal Society of London A} \textbf{249},
180-194 (1959). DOI: 10.1098/rspa.1959.0015.

{[}9{]} A. Pound and E. Poisson, ``Osculating Orbits in Schwarzschild
Spacetime, with an Application to Extreme Mass-Ratio Inspirals,''
\emph{Physical Review D} \textbf{77}, 044013 (2008). DOI:
10.1103/PhysRevD.77.044013.

{[}10{]} L. Blanchet, ``Post-Newtonian Theory for Gravitational Waves,''
\emph{Living Reviews in Relativity} \textbf{27}, 4 (2024). DOI:
10.1007/s41114-024-00050-z.

{[}11{]} L. Blanchet, G. Faye, and D. Trestini, ``Gravitational
Radiation Reaction for Compact Binary Systems at the Fourth-and-a-Half
Post-Newtonian Order,'' \emph{Classical and Quantum Gravity}
\textbf{42}, 065015 (2025). DOI: 10.1088/1361-6382/adac9d.

\begin{center}\rule{0.5\linewidth}{0.5pt}\end{center}

\section{Appendix A: Figures used in the main
text}\label{appendix-a-figures-used-in-the-main-text}

The figures of the revised main text are stored in

\begin{Shaded}
\begin{Highlighting}[]
\NormalTok{paper5\_revision\_v2\_figures\_tables\_20260926/}
\NormalTok{figures\_png\_preview/}
\end{Highlighting}
\end{Shaded}

\begin{enumerate}
\def\labelenumi{\arabic{enumi}.}
\tightlist
\item
  \texttt{paper5\_\allowbreak{}revision\_\allowbreak{}v2\_\allowbreak{}figures\_\allowbreak{}tables\_\allowbreak{}20260926/\allowbreak{}figures\_\allowbreak{}png\_\allowbreak{}preview/\allowbreak{}figure01.png}
  - the path of the state audit, 2 to 5 to 6
\item
  \texttt{08\_\allowbreak{}nu\_\allowbreak{}sixth\_\allowbreak{}state\_\allowbreak{}internalization\_\allowbreak{}control\_\allowbreak{}experiment\_\allowbreak{}20260925/\allowbreak{}figures/\allowbreak{}figure\_\allowbreak{}orbit\_\allowbreak{}c1\_\allowbreak{}paper4\_\allowbreak{}vs\_\allowbreak{}6state.png}
  - the Paper 4 C1 orbit and the orbit regenerated after internalizing
  the sixth state
\item
  \texttt{paper5\_\allowbreak{}revision\_\allowbreak{}v2\_\allowbreak{}figures\_\allowbreak{}tables\_\allowbreak{}20260926/\allowbreak{}figures\_\allowbreak{}png\_\allowbreak{}preview/\allowbreak{}figure02.png}
  - the scope of state closure, synchronous update and non-feedback
  readout
\item
  \texttt{paper5\_\allowbreak{}revision\_\allowbreak{}v2\_\allowbreak{}figures\_\allowbreak{}tables\_\allowbreak{}20260926/\allowbreak{}figures\_\allowbreak{}png\_\allowbreak{}preview/\allowbreak{}figure03.png}
  - 12-condition strict audit of scheme B
\item
  \texttt{paper5\_\allowbreak{}revision\_\allowbreak{}v2\_\allowbreak{}figures\_\allowbreak{}tables\_\allowbreak{}20260926/\allowbreak{}figures\_\allowbreak{}png\_\allowbreak{}preview/\allowbreak{}figure04.png}
  - convergence of the angular series for \(N=12,16,20,24\)
\item
  \texttt{paper5\_\allowbreak{}revision\_\allowbreak{}v2\_\allowbreak{}figures\_\allowbreak{}tables\_\allowbreak{}20260926/\allowbreak{}figures\_\allowbreak{}png\_\allowbreak{}preview/\allowbreak{}figure05.png}
  - convergence of the numerical integration at orders 4, 6 and 8
\item
  \texttt{paper5\_\allowbreak{}revision\_\allowbreak{}v2\_\allowbreak{}figures\_\allowbreak{}tables\_\allowbreak{}20260926/\allowbreak{}figures\_\allowbreak{}png\_\allowbreak{}preview/\allowbreak{}figure06.png}
  - comparison of the raising of the order over 9 conditions
\end{enumerate}

The same revision package includes the SVG originals, the table CSV
files, the source data, the figure/table regeneration program and the
reproducibility ZIP.

\begin{center}\rule{0.5\linewidth}{0.5pt}\end{center}

\section{Appendix B: Main files for
reproduction}\label{appendix-b-main-files-for-reproduction}

\begin{itemize}
\tightlist
\item
  \texttt{08\_\allowbreak{}nu\_\allowbreak{}sixth\_\allowbreak{}state\_\allowbreak{}internalization\_\allowbreak{}control\_\allowbreak{}experiment\_\allowbreak{}20260925/\allowbreak{}run\_\allowbreak{}nu\_\allowbreak{}internal\_\allowbreak{}state\_\allowbreak{}control\_\allowbreak{}v1.py}
\item
  \texttt{08\_\allowbreak{}nu\_\allowbreak{}sixth\_\allowbreak{}state\_\allowbreak{}internalization\_\allowbreak{}control\_\allowbreak{}experiment\_\allowbreak{}20260925/\allowbreak{}plot\_\allowbreak{}c1\_\allowbreak{}orbit\_\allowbreak{}6state\_\allowbreak{}audit\_\allowbreak{}v1.py}
\item
  \texttt{08\_\allowbreak{}nu\_\allowbreak{}sixth\_\allowbreak{}state\_\allowbreak{}internalization\_\allowbreak{}control\_\allowbreak{}experiment\_\allowbreak{}20260925/\allowbreak{}figures/\allowbreak{}figure\_\allowbreak{}orbit\_\allowbreak{}c1\_\allowbreak{}paper4\_\allowbreak{}vs\_\allowbreak{}6state.svg}
\item
  \texttt{08\_\allowbreak{}nu\_\allowbreak{}sixth\_\allowbreak{}state\_\allowbreak{}internalization\_\allowbreak{}control\_\allowbreak{}experiment\_\allowbreak{}20260925/\allowbreak{}figures/\allowbreak{}figure\_\allowbreak{}orbit\_\allowbreak{}c1\_\allowbreak{}paper4\_\allowbreak{}vs\_\allowbreak{}6state\_\allowbreak{}summary.txt}
\item
  \texttt{08\_\allowbreak{}nu\_\allowbreak{}sixth\_\allowbreak{}state\_\allowbreak{}internalization\_\allowbreak{}control\_\allowbreak{}experiment\_\allowbreak{}20260925/\allowbreak{}nu\_\allowbreak{}internal\_\allowbreak{}state\_\allowbreak{}control\_\allowbreak{}case\_\allowbreak{}summary.csv}
\item
  \texttt{08\_\allowbreak{}nu\_\allowbreak{}sixth\_\allowbreak{}state\_\allowbreak{}internalization\_\allowbreak{}control\_\allowbreak{}experiment\_\allowbreak{}20260925/\allowbreak{}nu\_\allowbreak{}internal\_\allowbreak{}state\_\allowbreak{}control\_\allowbreak{}summary.json}
\item
  \texttt{08\_\allowbreak{}nu\_\allowbreak{}sixth\_\allowbreak{}state\_\allowbreak{}internalization\_\allowbreak{}control\_\allowbreak{}experiment\_\allowbreak{}20260925/\allowbreak{}long\_\allowbreak{}run\_\allowbreak{}checks.json}
\item
  \texttt{08\_\allowbreak{}nu\_\allowbreak{}sixth\_\allowbreak{}state\_\allowbreak{}internalization\_\allowbreak{}control\_\allowbreak{}experiment\_\allowbreak{}20260925/\allowbreak{}saved\_\allowbreak{}C1\_\allowbreak{}check.json}
\item
  \texttt{paper5\_\allowbreak{}revision\_\allowbreak{}v2\_\allowbreak{}figures\_\allowbreak{}tables\_\allowbreak{}20260926/\allowbreak{}EXPERIMENT\_\allowbreak{}INDEX\_\allowbreak{}v2.csv}
\item
  \texttt{paper5\_\allowbreak{}revision\_\allowbreak{}v2\_\allowbreak{}figures\_\allowbreak{}tables\_\allowbreak{}20260926/\allowbreak{}FIGURE\_\allowbreak{}INDEX\_\allowbreak{}v2\_\allowbreak{}ja.md}
\item
  \texttt{paper5\_\allowbreak{}revision\_\allowbreak{}v2\_\allowbreak{}figures\_\allowbreak{}tables\_\allowbreak{}20260926/\allowbreak{}tables/\allowbreak{}table01\_\allowbreak{}state\_\allowbreak{}audit\_\allowbreak{}summary\_\allowbreak{}v2.csv}
\item
  \texttt{paper5\_\allowbreak{}revision\_\allowbreak{}v2\_\allowbreak{}figures\_\allowbreak{}tables\_\allowbreak{}20260926/\allowbreak{}tables/\allowbreak{}table02\_\allowbreak{}approximation\_\allowbreak{}and\_\allowbreak{}state\_\allowbreak{}scope\_\allowbreak{}v2.csv}
\item
  \texttt{paper5\_\allowbreak{}revision\_\allowbreak{}v2\_\allowbreak{}figures\_\allowbreak{}tables\_\allowbreak{}20260926/\allowbreak{}tables/\allowbreak{}table03\_\allowbreak{}key\_\allowbreak{}numerical\_\allowbreak{}results\_\allowbreak{}v2.csv}
\item
  \texttt{paper5\_\allowbreak{}revision\_\allowbreak{}v2\_\allowbreak{}figures\_\allowbreak{}tables\_\allowbreak{}20260926/\allowbreak{}scripts/\allowbreak{}generate\_\allowbreak{}paper5\_\allowbreak{}revision\_\allowbreak{}v2.py}
\item
  \texttt{paper5\_\allowbreak{}revision\_\allowbreak{}v2\_\allowbreak{}figures\_\allowbreak{}tables\_\allowbreak{}20260926/\allowbreak{}paper5\_\allowbreak{}revision\_\allowbreak{}v2\_\allowbreak{}figures\_\allowbreak{}tables\_\allowbreak{}20260926.zip}
\end{itemize}

The charged eight-state experiments (09 onwards) are not used as
evidence in this paper.

\end{document}
